\documentclass[
  doc,
  floatsintext,
  longtable,
  a4paper,
  nolmodern,
  notxfonts,
  notimes,
  12pt,
  colorlinks=true,linkcolor=blue,citecolor=blue,urlcolor=blue]{apa7}

\usepackage{amsmath}
\usepackage{amssymb}

\geometry{inner=1in, outer=1in}
\fancyhfoffset[LE,RO]{0cm}


\usepackage[bidi=default]{babel}
\babelprovide[main,import]{spanish}


% get rid of language-specific shorthands (see #6817):
\let\LanguageShortHands\languageshorthands
\def\languageshorthands#1{}

\RequirePackage{longtable}
\RequirePackage{threeparttablex}

\makeatletter
\renewcommand{\paragraph}{\@startsection{paragraph}{4}{\parindent}%
	{0\baselineskip \@plus 0.2ex \@minus 0.2ex}%
	{-.5em}%
	{\normalfont\normalsize\bfseries\typesectitle}}

\renewcommand{\subparagraph}[1]{\@startsection{subparagraph}{5}{0.5em}%
	{0\baselineskip \@plus 0.2ex \@minus 0.2ex}%
	{-\z@\relax}%
	{\normalfont\normalsize\bfseries\itshape\hspace{\parindent}{#1}\textit{\addperi}}{\relax}}
\makeatother




\usepackage{longtable, booktabs, multirow, multicol, colortbl, hhline, caption, array, float, xpatch}
\usepackage{subcaption}


\renewcommand\thesubfigure{\Alph{subfigure}}
\setcounter{topnumber}{2}
\setcounter{bottomnumber}{2}
\setcounter{totalnumber}{4}
\renewcommand{\topfraction}{0.85}
\renewcommand{\bottomfraction}{0.85}
\renewcommand{\textfraction}{0.15}
\renewcommand{\floatpagefraction}{0.7}

\usepackage{tcolorbox}
\tcbuselibrary{listings,theorems, breakable, skins}
\usepackage{fontawesome5}

\definecolor{quarto-callout-color}{HTML}{909090}
\definecolor{quarto-callout-note-color}{HTML}{0758E5}
\definecolor{quarto-callout-important-color}{HTML}{CC1914}
\definecolor{quarto-callout-warning-color}{HTML}{EB9113}
\definecolor{quarto-callout-tip-color}{HTML}{00A047}
\definecolor{quarto-callout-caution-color}{HTML}{FC5300}
\definecolor{quarto-callout-color-frame}{HTML}{ACACAC}
\definecolor{quarto-callout-note-color-frame}{HTML}{4582EC}
\definecolor{quarto-callout-important-color-frame}{HTML}{D9534F}
\definecolor{quarto-callout-warning-color-frame}{HTML}{F0AD4E}
\definecolor{quarto-callout-tip-color-frame}{HTML}{02B875}
\definecolor{quarto-callout-caution-color-frame}{HTML}{FD7E14}

%\newlength\Oldarrayrulewidth
%\newlength\Oldtabcolsep


\usepackage{hyperref}



\usepackage{color}
\usepackage{fancyvrb}
\newcommand{\VerbBar}{|}
\newcommand{\VERB}{\Verb[commandchars=\\\{\}]}
\DefineVerbatimEnvironment{Highlighting}{Verbatim}{commandchars=\\\{\}}
% Add ',fontsize=\small' for more characters per line
\usepackage{framed}
\definecolor{shadecolor}{RGB}{241,243,245}
\newenvironment{Shaded}{\begin{snugshade}}{\end{snugshade}}
\newcommand{\AlertTok}[1]{\textcolor[rgb]{0.68,0.00,0.00}{#1}}
\newcommand{\AnnotationTok}[1]{\textcolor[rgb]{0.37,0.37,0.37}{#1}}
\newcommand{\AttributeTok}[1]{\textcolor[rgb]{0.40,0.45,0.13}{#1}}
\newcommand{\BaseNTok}[1]{\textcolor[rgb]{0.68,0.00,0.00}{#1}}
\newcommand{\BuiltInTok}[1]{\textcolor[rgb]{0.00,0.23,0.31}{#1}}
\newcommand{\CharTok}[1]{\textcolor[rgb]{0.13,0.47,0.30}{#1}}
\newcommand{\CommentTok}[1]{\textcolor[rgb]{0.37,0.37,0.37}{#1}}
\newcommand{\CommentVarTok}[1]{\textcolor[rgb]{0.37,0.37,0.37}{\textit{#1}}}
\newcommand{\ConstantTok}[1]{\textcolor[rgb]{0.56,0.35,0.01}{#1}}
\newcommand{\ControlFlowTok}[1]{\textcolor[rgb]{0.00,0.23,0.31}{\textbf{#1}}}
\newcommand{\DataTypeTok}[1]{\textcolor[rgb]{0.68,0.00,0.00}{#1}}
\newcommand{\DecValTok}[1]{\textcolor[rgb]{0.68,0.00,0.00}{#1}}
\newcommand{\DocumentationTok}[1]{\textcolor[rgb]{0.37,0.37,0.37}{\textit{#1}}}
\newcommand{\ErrorTok}[1]{\textcolor[rgb]{0.68,0.00,0.00}{#1}}
\newcommand{\ExtensionTok}[1]{\textcolor[rgb]{0.00,0.23,0.31}{#1}}
\newcommand{\FloatTok}[1]{\textcolor[rgb]{0.68,0.00,0.00}{#1}}
\newcommand{\FunctionTok}[1]{\textcolor[rgb]{0.28,0.35,0.67}{#1}}
\newcommand{\ImportTok}[1]{\textcolor[rgb]{0.00,0.46,0.62}{#1}}
\newcommand{\InformationTok}[1]{\textcolor[rgb]{0.37,0.37,0.37}{#1}}
\newcommand{\KeywordTok}[1]{\textcolor[rgb]{0.00,0.23,0.31}{\textbf{#1}}}
\newcommand{\NormalTok}[1]{\textcolor[rgb]{0.00,0.23,0.31}{#1}}
\newcommand{\OperatorTok}[1]{\textcolor[rgb]{0.37,0.37,0.37}{#1}}
\newcommand{\OtherTok}[1]{\textcolor[rgb]{0.00,0.23,0.31}{#1}}
\newcommand{\PreprocessorTok}[1]{\textcolor[rgb]{0.68,0.00,0.00}{#1}}
\newcommand{\RegionMarkerTok}[1]{\textcolor[rgb]{0.00,0.23,0.31}{#1}}
\newcommand{\SpecialCharTok}[1]{\textcolor[rgb]{0.37,0.37,0.37}{#1}}
\newcommand{\SpecialStringTok}[1]{\textcolor[rgb]{0.13,0.47,0.30}{#1}}
\newcommand{\StringTok}[1]{\textcolor[rgb]{0.13,0.47,0.30}{#1}}
\newcommand{\VariableTok}[1]{\textcolor[rgb]{0.07,0.07,0.07}{#1}}
\newcommand{\VerbatimStringTok}[1]{\textcolor[rgb]{0.13,0.47,0.30}{#1}}
\newcommand{\WarningTok}[1]{\textcolor[rgb]{0.37,0.37,0.37}{\textit{#1}}}

\providecommand{\tightlist}{%
  \setlength{\itemsep}{0pt}\setlength{\parskip}{0pt}}
\usepackage{longtable,booktabs,array}
\usepackage{calc} % for calculating minipage widths
% Correct order of tables after \paragraph or \subparagraph
\usepackage{etoolbox}
\makeatletter
\patchcmd\longtable{\par}{\if@noskipsec\mbox{}\fi\par}{}{}
\makeatother
% Allow footnotes in longtable head/foot
\IfFileExists{footnotehyper.sty}{\usepackage{footnotehyper}}{\usepackage{footnote}}
\makesavenoteenv{longtable}

\usepackage{graphicx}
\makeatletter
\newsavebox\pandoc@box
\newcommand*\pandocbounded[1]{% scales image to fit in text height/width
  \sbox\pandoc@box{#1}%
  \Gscale@div\@tempa{\textheight}{\dimexpr\ht\pandoc@box+\dp\pandoc@box\relax}%
  \Gscale@div\@tempb{\linewidth}{\wd\pandoc@box}%
  \ifdim\@tempb\p@<\@tempa\p@\let\@tempa\@tempb\fi% select the smaller of both
  \ifdim\@tempa\p@<\p@\scalebox{\@tempa}{\usebox\pandoc@box}%
  \else\usebox{\pandoc@box}%
  \fi%
}
% Set default figure placement to htbp
\def\fps@figure{htbp}
\makeatother







\usepackage{newtx}

\defaultfontfeatures{Scale=MatchLowercase}
\defaultfontfeatures[\rmfamily]{Ligatures=TeX,Scale=1}





\title{Guía Completa de Kitty Terminal}


\shorttitle{Guía Completa de Kitty Terminal}


\usepackage{etoolbox}



\ccoppy{\textcopyright~2025}



\author{Edison Achalma}



\affiliation{
{Escuela Profesional de Economía, Universidad Nacional de San Cristóbal
de Huamanga}}




\leftheader{Achalma}

\date{2025-12-31}


\abstract{Kitty es un emulador de terminal moderno, rápido y altamente
personalizable que aprovecha la aceleración por GPU para ofrecer un
rendimiento excepcional. Esta guía completa aborda desde la instalación
en diferentes sistemas operativos hasta la configuración avanzada y
optimización profunda de Kitty. Se exploran en detalle; arquitectura y
principios de diseño, jerarquía de ventanas/pestañas, configuración de
kitty.conf, manejo de fuentes con ligaduras, gestión avanzada de colores
y temas, layouts nativos (tall, fat, grid, splits, stack), kittens
(extensiones), control remoto, integración con shell, marcas, sesiones,
optimización de rendimiento y solución de problemas frecuentes. Dirigida
tanto a usuarios que buscan reemplazar terminales tradicionales por una
opción más veloz y moderna, como a usuarios avanzados que desean
exprimir al máximo las capacidades de personalización y automatización
que ofrece Kitty en entornos de desarrollo y administración de
sistemas. }

\keywords{Kitty Terminal, Linux, Linea de comandos}

\authornote{\par{\addORCIDlink{Edison Achalma}{0000-0001-6996-3364}} 
\par{ }
\par{   El autor no tiene conflictos de interés que revelar.    Los
roles de autor se clasificaron utilizando la taxonomía de roles de
colaborador (CRediT; https://credit.niso.org/) de la siguiente
manera:  Edison Achalma:   conceptualización, redacción}
\par{La correspondencia relativa a este artículo debe dirigirse a Edison
Achalma, Email: \href{mailto:elmer.achalma.09@unsch.edu.pe}{elmer.achalma.09@unsch.edu.pe}}
}

\makeatletter
\let\endoldlt\endlongtable
\def\endlongtable{
\hline
\endoldlt
}
\makeatother

\urlstyle{same}



\makeatletter
\@ifpackageloaded{caption}{}{\usepackage{caption}}
\AtBeginDocument{%
\ifdefined\contentsname
  \renewcommand*\contentsname{Tabla de contenidos}
\else
  \newcommand\contentsname{Tabla de contenidos}
\fi
\ifdefined\listfigurename
  \renewcommand*\listfigurename{Lista de Figuras}
\else
  \newcommand\listfigurename{Lista de Figuras}
\fi
\ifdefined\listtablename
  \renewcommand*\listtablename{Lista de Tablas}
\else
  \newcommand\listtablename{Lista de Tablas}
\fi
\ifdefined\figurename
  \renewcommand*\figurename{Figura}
\else
  \newcommand\figurename{Figura}
\fi
\ifdefined\tablename
  \renewcommand*\tablename{Tabla}
\else
  \newcommand\tablename{Tabla}
\fi
}
\@ifpackageloaded{float}{}{\usepackage{float}}
\floatstyle{ruled}
\@ifundefined{c@chapter}{\newfloat{codelisting}{h}{lop}}{\newfloat{codelisting}{h}{lop}[chapter]}
\floatname{codelisting}{Listing}
\newcommand*\listoflistings{\listof{codelisting}{Lista de Listings}}
\makeatother
\makeatletter
\makeatother
\makeatletter
\@ifpackageloaded{caption}{}{\usepackage{caption}}
\@ifpackageloaded{subcaption}{}{\usepackage{subcaption}}
\makeatother
\makeatletter
\@ifpackageloaded{fontawesome5}{}{\usepackage{fontawesome5}}
\makeatother

% From https://tex.stackexchange.com/a/645996/211326
%%% apa7 doesn't want to add appendix section titles in the toc
%%% let's make it do it
\makeatletter
\xpatchcmd{\appendix}
  {\par}
  {\addcontentsline{toc}{section}{\@currentlabelname}\par}
  {}{}
\makeatother

%% Disable longtable counter
%% https://tex.stackexchange.com/a/248395/211326

\usepackage{etoolbox}

\makeatletter
\patchcmd{\LT@caption}
  {\bgroup}
  {\bgroup\global\LTpatch@captiontrue}
  {}{}
\patchcmd{\longtable}
  {\par}
  {\par\global\LTpatch@captionfalse}
  {}{}
\apptocmd{\endlongtable}
  {\ifLTpatch@caption\else\addtocounter{table}{-1}\fi}
  {}{}
\newif\ifLTpatch@caption
\makeatother

\begin{document}

\maketitle


\hypertarget{toc}{}
\tableofcontents
\newpage
\section[Introduction]{Guía Completa de Kitty Terminal}

\setcounter{secnumdepth}{3}

\setlength\LTleft{0pt}




\textbf{¿Qué es Kitty?}

Kitty es un emulador de terminal moderno, rápido y rico en
características, diseñado desde cero para aprovechar las GPUs modernas.
Creado por Kovid Goyal, es conocido por:

\textbf{Características Principales:}

\begin{itemize}
\tightlist
\item
  \textbf{Renderizado GPU acelerado}: Extremadamente rápido
\item
  \textbf{Soporte Unicode completo}: Incluyendo emojis y ligaduras
\item
  \textbf{True color}: 24-bit RGB
\item
  \textbf{Ligaduras de fuentes}: Soporte nativo
\item
  \textbf{Imágenes en terminal}: Protocolo propio de alta calidad
\item
  \textbf{Multiplexación}: Tabs y splits nativos
\item
  \textbf{Extensible}: Sistema de kittens (plugins)
\item
  \textbf{Control remoto}: API completa
\item
  \textbf{Configuración simple}: Archivo único y legible
\end{itemize}

\textbf{Por Qué Usar Kitty}

\textbf{Ventajas sobre otros terminales:}

\begin{longtable}[]{@{}
  >{\raggedright\arraybackslash}p{(\linewidth - 8\tabcolsep) * \real{0.2809}}
  >{\raggedright\arraybackslash}p{(\linewidth - 8\tabcolsep) * \real{0.1798}}
  >{\raggedright\arraybackslash}p{(\linewidth - 8\tabcolsep) * \real{0.1798}}
  >{\raggedright\arraybackslash}p{(\linewidth - 8\tabcolsep) * \real{0.1798}}
  >{\raggedright\arraybackslash}p{(\linewidth - 8\tabcolsep) * \real{0.1798}}@{}}
\toprule\noalign{}
\begin{minipage}[b]{\linewidth}\raggedright
Característica
\end{minipage} & \begin{minipage}[b]{\linewidth}\raggedright
Kitty
\end{minipage} & \begin{minipage}[b]{\linewidth}\raggedright
Alacritty
\end{minipage} & \begin{minipage}[b]{\linewidth}\raggedright
iTerm2
\end{minipage} & \begin{minipage}[b]{\linewidth}\raggedright
GNOME Terminal
\end{minipage} \\
\midrule\noalign{}
\endhead
\bottomrule\noalign{}
\endlastfoot
Renderizado por GPU & Sí & Sí & Sí & No \\
Pestañas nativas & Sí & No & Sí & Sí \\
División de paneles & Sí & No & Sí & No \\
Soporte para imágenes & Sí & No & Sí & No \\
Ligaduras (ligatures) & Sí & Sí & Sí & Sí \\
Sistema de extensiones & Sí (kittens) & No & Sí & No \\
Control remoto (SSH) & Sí & No & Sí & No \\
Multiplataforma & Sí (Linux/macOS/Windows) & Sí (Linux/macOS/Windows) &
No (solo macOS) & No (solo Linux) \\
\end{longtable}

\section{Instalación}\label{instalaciuxf3n}

\subsection{Requisitos del Sistema}\label{requisitos-del-sistema}

\textbf{Mínimos:}

\begin{itemize}
\tightlist
\item
  OpenGL 3.3+
\item
  Linux, macOS, o Windows (WSL)
\item
  Python 3.6+
\end{itemize}

\textbf{Recomendados:}

\begin{itemize}
\tightlist
\item
  GPU moderna con drivers actualizados
\item
  Monitor de alta resolución (4K/Retina)
\item
  Fuente con ligaduras (FiraCode, JetBrains Mono)
\end{itemize}

\subsection{Linux}\label{linux}

\textbf{Ubuntu/Debian:}

\begin{Shaded}
\begin{Highlighting}[]
\FunctionTok{sudo}\NormalTok{ apt install kitty}
\end{Highlighting}
\end{Shaded}

\textbf{Arch Linux:}

\begin{Shaded}
\begin{Highlighting}[]
\FunctionTok{sudo}\NormalTok{ pacman }\AttributeTok{{-}S}\NormalTok{ kitty}
\end{Highlighting}
\end{Shaded}

\textbf{Fedora:}

\begin{Shaded}
\begin{Highlighting}[]
\FunctionTok{sudo}\NormalTok{ dnf install kitty}
\end{Highlighting}
\end{Shaded}

\textbf{Desde binario oficial (recomendado):}

\begin{Shaded}
\begin{Highlighting}[]
\ExtensionTok{curl} \AttributeTok{{-}L}\NormalTok{ https://sw.kovidgoyal.net/kitty/installer.sh }\KeywordTok{|} \FunctionTok{sh}\NormalTok{ /dev/stdin}
\end{Highlighting}
\end{Shaded}

\subsection{macOS}\label{macos}

\textbf{Homebrew:}

\begin{Shaded}
\begin{Highlighting}[]
\ExtensionTok{brew}\NormalTok{ install }\AttributeTok{{-}{-}cask}\NormalTok{ kitty}
\end{Highlighting}
\end{Shaded}

\textbf{Desde DMG:}

Descargar desde: https://github.com/kovidgoyal/kitty/releases

\subsection{Windows (WSL)}\label{windows-wsl}

\begin{Shaded}
\begin{Highlighting}[]
\ExtensionTok{curl} \AttributeTok{{-}L}\NormalTok{ https://sw.kovidgoyal.net/kitty/installer.sh }\KeywordTok{|} \FunctionTok{sh}\NormalTok{ /dev/stdin}
\end{Highlighting}
\end{Shaded}

\subsection{Verificar Instalación}\label{verificar-instalaciuxf3n}

\begin{Shaded}
\begin{Highlighting}[]
\ExtensionTok{kitty} \AttributeTok{{-}{-}version}
\end{Highlighting}
\end{Shaded}

\subsection{Instalación de Dependencias
Opcionales}\label{instalaciuxf3n-de-dependencias-opcionales}

\begin{Shaded}
\begin{Highlighting}[]
\CommentTok{\# Herramientas útiles para kittens}
\FunctionTok{sudo}\NormalTok{ apt install imagemagick          }\CommentTok{\# Procesamiento de imágenes}
\FunctionTok{sudo}\NormalTok{ apt install bat                  }\CommentTok{\# Mejor syntax highlighting}
\FunctionTok{sudo}\NormalTok{ apt install fzf                  }\CommentTok{\# Fuzzy finder}
\FunctionTok{sudo}\NormalTok{ apt install ripgrep              }\CommentTok{\# Búsqueda rápida}
\end{Highlighting}
\end{Shaded}

\section{Filosofía y Diseño}\label{filosofuxeda-y-diseuxf1o}

\subsection{Principios de Diseño}\label{principios-de-diseuxf1o}

\begin{enumerate}
\def\labelenumi{\arabic{enumi}.}
\tightlist
\item
  \textbf{Velocidad}: GPU rendering para máximo rendimiento
\item
  \textbf{Simplicidad}: Configuración en un solo archivo
\item
  \textbf{Extensibilidad}: Kittens para añadir funcionalidad
\item
  \textbf{Estándares}: Soporte completo de especificaciones modernas
\item
  \textbf{Teclado primero}: Control total desde teclado
\end{enumerate}

\subsection{Arquitectura}\label{arquitectura}

\textbf{Componentes principales:}

\begin{itemize}
\tightlist
\item
  \textbf{Core (C)}: Renderizado GPU, performance crítico
\item
  \textbf{UI (Python)}: Interfaz, configuración, extensibilidad
\item
  \textbf{Kittens (Python/Go)}: Programas auxiliares
\end{itemize}

\textbf{Stack tecnológico:}

\begin{itemize}
\tightlist
\item
  OpenGL para renderizado
\item
  HarfBuzz para text shaping
\item
  FontConfig para gestión de fuentes
\item
  libunistring para Unicode
\end{itemize}

\section{Conceptos Básicos}\label{conceptos-buxe1sicos}

\subsection{Jerarquía: OS Windows \textgreater{} Tabs \textgreater{}
Windows}\label{jerarquuxeda-os-windows-tabs-windows}

\begin{verbatim}
┌─────────────────────────────────────┐
│ OS Window (ventana del sistema)     │
│ ┌───────────────────────────────┐   │
│ │ Tab 1    Tab 2    Tab 3       │   │
│ ├───────────────────────────────┤   │
│ │ ┌──────────┬──────────┐       │   │
│ │ │ Window 1 │ Window 2 │       │   │
│ │ │          │          │       │   │
│ │ └──────────┴──────────┘       │   │
│ └───────────────────────────────┘   │
└─────────────────────────────────────┘
\end{verbatim}

\textbf{OS Window}: Ventana del sistema operativo \textbf{Tab}: Pestaña
dentro de un OS Window \textbf{Window}: Panel dentro de un Tab (como
splits en tmux)

\subsection{Primer Uso}\label{primer-uso}

\textbf{Al abrir Kitty por primera vez:}

\begin{Shaded}
\begin{Highlighting}[]
\CommentTok{\# Abrir Kitty}
\ExtensionTok{kitty}

\CommentTok{\# Ver ayuda}
\ExtensionTok{Ctrl+Shift+F1}

\CommentTok{\# Editar configuración}
\ExtensionTok{Ctrl+Shift+F2}

\CommentTok{\# Recargar configuración}
\ExtensionTok{Ctrl+Shift+F5}
\end{Highlighting}
\end{Shaded}

\section{Configuración (kitty.conf)}\label{configuraciuxf3n-kitty.conf}

\subsection{Ubicación del Archivo}\label{ubicaciuxf3n-del-archivo}

\begin{verbatim}
Linux/macOS: ~/.config/kitty/kitty.conf
Windows: %LOCALAPPDATA%\kitty\kitty.conf
\end{verbatim}

\subsection{Generar Configuración Por
Defecto}\label{generar-configuraciuxf3n-por-defecto}

\begin{Shaded}
\begin{Highlighting}[]
\CommentTok{\# Generar con comentarios}
\ExtensionTok{kitty}\NormalTok{ +runpy }\StringTok{\textquotesingle{}from kitty.config import *; print(commented\_out\_default\_config())\textquotesingle{}}

\CommentTok{\# O abrir para editar}
\ExtensionTok{Ctrl+Shift+F2}
\end{Highlighting}
\end{Shaded}

\subsection{Estructura de
Configuración}\label{estructura-de-configuraciuxf3n}

\textbf{Secciones principales:}

\begin{enumerate}
\def\labelenumi{\arabic{enumi}.}
\tightlist
\item
  Fuentes
\item
  Cursor
\item
  Scrollback
\item
  Mouse
\item
  Terminal bell
\item
  Window layout
\item
  Tab bar
\item
  Color scheme
\item
  Advanced
\item
  OS tweaks
\item
  Keyboard shortcuts
\end{enumerate}

\subsection{Configuración Básica
Recomendada}\label{configuraciuxf3n-buxe1sica-recomendada}

\begin{Shaded}
\begin{Highlighting}[]
\CommentTok{\# ==========================================}
\CommentTok{\# CONFIGURACIÓN BÁSICA DE KITTY}
\CommentTok{\# ==========================================}

\CommentTok{\# Fuentes}
\DataTypeTok{font\_family}      \DataTypeTok{JetBrains} \DataTypeTok{Mono}
\DataTypeTok{bold\_font}        \DataTypeTok{auto}
\DataTypeTok{italic\_font}      \DataTypeTok{auto}
\DataTypeTok{bold\_italic\_font} \DataTypeTok{auto}
\DataTypeTok{font\_size}        \DataTypeTok{12.0}

\CommentTok{\# Cursor}
\DataTypeTok{cursor\_shape} \DataTypeTok{block}
\DataTypeTok{cursor\_blink\_interval} \DataTypeTok{0}
\DataTypeTok{cursor\_stop\_blinking\_after} \DataTypeTok{15.0}

\CommentTok{\# Scrollback}
\DataTypeTok{scrollback\_lines} \DataTypeTok{10000}
\DataTypeTok{scrollback\_pager} \DataTypeTok{less} \DataTypeTok{{-}{-}chop{-}long{-}lines} \DataTypeTok{{-}{-}RAW{-}CONTROL{-}CHARS} \ErrorTok{+}\DataTypeTok{INPUT\_LINE\_NUMBER}

\CommentTok{\# Mouse}
\DataTypeTok{mouse\_hide\_wait} \DataTypeTok{3.0}
\DataTypeTok{url\_color} \CommentTok{\#0087bd}
\DataTypeTok{url\_style} \DataTypeTok{curly}
\DataTypeTok{open\_url\_with} \DataTypeTok{default}
\DataTypeTok{detect\_urls} \DataTypeTok{yes}
\DataTypeTok{copy\_on\_select} \DataTypeTok{no}

\CommentTok{\# Ventanas}
\DataTypeTok{remember\_window\_size}  \DataTypeTok{yes}
\DataTypeTok{initial\_window\_width}  \DataTypeTok{640}
\DataTypeTok{initial\_window\_height} \DataTypeTok{400}
\DataTypeTok{window\_border\_width} \DataTypeTok{0.5pt}
\DataTypeTok{window\_margin\_width} \DataTypeTok{0}
\DataTypeTok{window\_padding\_width} \DataTypeTok{0}
\DataTypeTok{active\_border\_color} \CommentTok{\#00ff00}
\DataTypeTok{inactive\_border\_color} \CommentTok{\#cccccc}

\CommentTok{\# Pestañas}
\DataTypeTok{tab\_bar\_edge} \DataTypeTok{bottom}
\DataTypeTok{tab\_bar\_style} \DataTypeTok{fade}
\DataTypeTok{tab\_bar\_min\_tabs} \DataTypeTok{2}
\DataTypeTok{tab\_title\_template} \DataTypeTok{"\{index\}: \{title\}"}

\CommentTok{\# Colores (Dracula theme)}
\DataTypeTok{foreground} \CommentTok{\#f8f8f2}
\DataTypeTok{background} \CommentTok{\#282a36}
\DataTypeTok{selection\_foreground} \CommentTok{\#ffffff}
\DataTypeTok{selection\_background} \CommentTok{\#44475a}

\CommentTok{\# Negro}
\DataTypeTok{color0}  \CommentTok{\#21222c}
\DataTypeTok{color8}  \CommentTok{\#6272a4}

\CommentTok{\# Rojo}
\DataTypeTok{color1}  \CommentTok{\#ff5555}
\DataTypeTok{color9}  \CommentTok{\#ff6e6e}

\CommentTok{\# Verde}
\DataTypeTok{color2}  \CommentTok{\#50fa7b}
\DataTypeTok{color10} \CommentTok{\#69ff94}

\CommentTok{\# Amarillo}
\DataTypeTok{color3}  \CommentTok{\#f1fa8c}
\DataTypeTok{color11} \CommentTok{\#ffffa5}

\CommentTok{\# Azul}
\DataTypeTok{color4}  \CommentTok{\#bd93f9}
\DataTypeTok{color12} \CommentTok{\#d6acff}

\CommentTok{\# Magenta}
\DataTypeTok{color5}  \CommentTok{\#ff79c6}
\DataTypeTok{color13} \CommentTok{\#ff92df}

\CommentTok{\# Cyan}
\DataTypeTok{color6}  \CommentTok{\#8be9fd}
\DataTypeTok{color14} \CommentTok{\#a4ffff}

\CommentTok{\# Blanco}
\DataTypeTok{color7}  \CommentTok{\#f8f8f2}
\DataTypeTok{color15} \CommentTok{\#ffffff}

\CommentTok{\# Rendimiento}
\DataTypeTok{repaint\_delay} \DataTypeTok{10}
\DataTypeTok{input\_delay} \DataTypeTok{3}
\DataTypeTok{sync\_to\_monitor} \DataTypeTok{yes}

\CommentTok{\# Bell}
\DataTypeTok{enable\_audio\_bell} \DataTypeTok{no}
\DataTypeTok{visual\_bell\_duration} \DataTypeTok{0.0}
\DataTypeTok{window\_alert\_on\_bell} \DataTypeTok{yes}
\DataTypeTok{bell\_on\_tab} \DataTypeTok{"🔔 "}

\CommentTok{\# Advanced}
\DataTypeTok{shell} \DataTypeTok{.}
\DataTypeTok{editor} \DataTypeTok{.}
\DataTypeTok{allow\_remote\_control} \DataTypeTok{yes}
\DataTypeTok{listen\_on} \DataTypeTok{unix}\ErrorTok{:/}\DataTypeTok{tmp}\ErrorTok{/}\DataTypeTok{kitty}
\DataTypeTok{clipboard\_control} \DataTypeTok{write{-}clipboard} \DataTypeTok{write{-}primary} \DataTypeTok{read{-}clipboard{-}ask} \DataTypeTok{read{-}primary{-}ask}

\CommentTok{\# OS Específico (macOS)}
\DataTypeTok{macos\_titlebar\_color} \DataTypeTok{system}
\DataTypeTok{macos\_option\_as\_alt} \DataTypeTok{no}
\DataTypeTok{macos\_hide\_from\_tasks} \DataTypeTok{no}
\DataTypeTok{macos\_quit\_when\_last\_window\_closed} \DataTypeTok{no}

\CommentTok{\# Layouts habilitados}
\DataTypeTok{enabled\_layouts} \DataTypeTok{tall}\ErrorTok{:}\DataTypeTok{bias}\OperatorTok{=}\DecValTok{50}\ErrorTok{;}\DataTypeTok{full\_size}\OperatorTok{=}\DecValTok{1}\ErrorTok{,}\DataTypeTok{stack}\ErrorTok{,}\DataTypeTok{fat}\ErrorTok{:}\DataTypeTok{bias}\OperatorTok{=}\DecValTok{50}\ErrorTok{;}\DataTypeTok{full\_size}\OperatorTok{=}\DecValTok{1}\ErrorTok{,}\DataTypeTok{grid}\ErrorTok{,}\DataTypeTok{splits}

\CommentTok{\# ==========================================}
\CommentTok{\# MAPEOS DE TECLADO}
\CommentTok{\# ==========================================}

\CommentTok{\# Modificador principal}
\DataTypeTok{kitty\_mod} \DataTypeTok{ctrl}\ErrorTok{+}\DataTypeTok{shift}

\CommentTok{\# Clipboard}
\DataTypeTok{map} \DataTypeTok{kitty\_mod}\ErrorTok{+}\DataTypeTok{c} \DataTypeTok{copy\_to\_clipboard}
\DataTypeTok{map} \DataTypeTok{kitty\_mod}\ErrorTok{+}\DataTypeTok{v} \DataTypeTok{paste\_from\_clipboard}
\DataTypeTok{map} \DataTypeTok{kitty\_mod}\ErrorTok{+}\DataTypeTok{s} \DataTypeTok{paste\_from\_selection}

\CommentTok{\# Scrolling}
\DataTypeTok{map} \DataTypeTok{kitty\_mod}\ErrorTok{+}\DataTypeTok{up} \DataTypeTok{scroll\_line\_up}
\DataTypeTok{map} \DataTypeTok{kitty\_mod}\ErrorTok{+}\DataTypeTok{down} \DataTypeTok{scroll\_line\_down}
\DataTypeTok{map} \DataTypeTok{kitty\_mod}\ErrorTok{+}\DataTypeTok{page\_up} \DataTypeTok{scroll\_page\_up}
\DataTypeTok{map} \DataTypeTok{kitty\_mod}\ErrorTok{+}\DataTypeTok{page\_down} \DataTypeTok{scroll\_page\_down}
\DataTypeTok{map} \DataTypeTok{kitty\_mod}\ErrorTok{+}\DataTypeTok{home} \DataTypeTok{scroll\_home}
\DataTypeTok{map} \DataTypeTok{kitty\_mod}\ErrorTok{+}\DataTypeTok{end} \DataTypeTok{scroll\_end}
\DataTypeTok{map} \DataTypeTok{kitty\_mod}\ErrorTok{+}\DataTypeTok{h} \DataTypeTok{show\_scrollback}
\DataTypeTok{map} \DataTypeTok{kitty\_mod}\ErrorTok{+}\DataTypeTok{g} \DataTypeTok{show\_last\_command\_output}

\CommentTok{\# Ventanas}
\DataTypeTok{map} \DataTypeTok{kitty\_mod}\ErrorTok{+}\DataTypeTok{enter} \DataTypeTok{new\_window}
\DataTypeTok{map} \DataTypeTok{kitty\_mod}\ErrorTok{+}\DataTypeTok{n} \DataTypeTok{new\_os\_window}
\DataTypeTok{map} \DataTypeTok{kitty\_mod}\ErrorTok{+}\DataTypeTok{w} \DataTypeTok{close\_window}
\DataTypeTok{map} \DataTypeTok{kitty\_mod}\ErrorTok{+]} \DataTypeTok{next\_window}
\DataTypeTok{map} \DataTypeTok{kitty\_mod}\ErrorTok{+}\KeywordTok{[} \KeywordTok{previous\_window}
\DataTypeTok{map} \DataTypeTok{kitty\_mod}\ErrorTok{+}\DataTypeTok{f} \DataTypeTok{move\_window\_forward}
\DataTypeTok{map} \DataTypeTok{kitty\_mod}\ErrorTok{+}\DataTypeTok{b} \DataTypeTok{move\_window\_backward}
\DataTypeTok{map} \DataTypeTok{kitty\_mod}\ErrorTok{+\textasciigrave{}} \DataTypeTok{move\_window\_to\_top}
\DataTypeTok{map} \DataTypeTok{kitty\_mod}\ErrorTok{+}\DataTypeTok{r} \DataTypeTok{start\_resizing\_window}
\DataTypeTok{map} \DataTypeTok{kitty\_mod}\ErrorTok{+}\DataTypeTok{1} \DataTypeTok{first\_window}
\DataTypeTok{map} \DataTypeTok{kitty\_mod}\ErrorTok{+}\DataTypeTok{2} \DataTypeTok{second\_window}
\DataTypeTok{map} \DataTypeTok{kitty\_mod}\ErrorTok{+}\DataTypeTok{3} \DataTypeTok{third\_window}
\DataTypeTok{map} \DataTypeTok{kitty\_mod}\ErrorTok{+}\DataTypeTok{4} \DataTypeTok{fourth\_window}
\DataTypeTok{map} \DataTypeTok{kitty\_mod}\ErrorTok{+}\DataTypeTok{5} \DataTypeTok{fifth\_window}

\CommentTok{\# Tabs}
\DataTypeTok{map} \DataTypeTok{kitty\_mod}\ErrorTok{+}\DataTypeTok{t} \DataTypeTok{new\_tab}
\DataTypeTok{map} \DataTypeTok{kitty\_mod}\ErrorTok{+}\DataTypeTok{q} \DataTypeTok{close\_tab}
\DataTypeTok{map} \DataTypeTok{kitty\_mod}\ErrorTok{+}\DataTypeTok{right} \DataTypeTok{next\_tab}
\DataTypeTok{map} \DataTypeTok{kitty\_mod}\ErrorTok{+}\DataTypeTok{left} \DataTypeTok{previous\_tab}
\DataTypeTok{map} \DataTypeTok{kitty\_mod}\ErrorTok{+}\DataTypeTok{.} \DataTypeTok{move\_tab\_forward}
\DataTypeTok{map} \DataTypeTok{kitty\_mod}\ErrorTok{+,} \DataTypeTok{move\_tab\_backward}
\DataTypeTok{map} \DataTypeTok{kitty\_mod}\ErrorTok{+}\DataTypeTok{alt}\ErrorTok{+}\DataTypeTok{t} \DataTypeTok{set\_tab\_title}

\CommentTok{\# Layouts}
\DataTypeTok{map} \DataTypeTok{kitty\_mod}\ErrorTok{+}\DataTypeTok{l} \DataTypeTok{next\_layout}

\CommentTok{\# Fuentes}
\DataTypeTok{map} \DataTypeTok{kitty\_mod}\ErrorTok{+}\DataTypeTok{equal} \DataTypeTok{change\_font\_size} \DataTypeTok{all} \ErrorTok{+}\DataTypeTok{2.0}
\DataTypeTok{map} \DataTypeTok{kitty\_mod}\ErrorTok{+}\DataTypeTok{minus} \DataTypeTok{change\_font\_size} \DataTypeTok{all} \DataTypeTok{{-}2.0}
\DataTypeTok{map} \DataTypeTok{kitty\_mod}\ErrorTok{+}\DataTypeTok{backspace} \DataTypeTok{change\_font\_size} \DataTypeTok{all} \DataTypeTok{0}

\CommentTok{\# Misc}
\DataTypeTok{map} \DataTypeTok{kitty\_mod}\ErrorTok{+}\DataTypeTok{f1} \DataTypeTok{show\_kitty\_doc} \DataTypeTok{overview}
\DataTypeTok{map} \DataTypeTok{kitty\_mod}\ErrorTok{+}\DataTypeTok{f2} \DataTypeTok{edit\_config\_file}
\DataTypeTok{map} \DataTypeTok{kitty\_mod}\ErrorTok{+}\DataTypeTok{f5} \DataTypeTok{load\_config\_file}
\DataTypeTok{map} \DataTypeTok{kitty\_mod}\ErrorTok{+}\DataTypeTok{f6} \DataTypeTok{debug\_config}
\DataTypeTok{map} \DataTypeTok{kitty\_mod}\ErrorTok{+}\DataTypeTok{f11} \DataTypeTok{toggle\_fullscreen}
\DataTypeTok{map} \DataTypeTok{kitty\_mod}\ErrorTok{+}\DataTypeTok{f10} \DataTypeTok{toggle\_maximized}
\DataTypeTok{map} \DataTypeTok{kitty\_mod}\ErrorTok{+}\DataTypeTok{u} \DataTypeTok{kitten} \DataTypeTok{unicode\_input}
\DataTypeTok{map} \DataTypeTok{kitty\_mod}\ErrorTok{+}\DataTypeTok{e} \DataTypeTok{open\_url\_with\_hints}
\DataTypeTok{map} \DataTypeTok{kitty\_mod}\ErrorTok{+}\DataTypeTok{escape} \DataTypeTok{kitty\_shell} \DataTypeTok{window}

\CommentTok{\# Opacidad de fondo}
\DataTypeTok{map} \DataTypeTok{kitty\_mod}\ErrorTok{+}\DataTypeTok{a}\ErrorTok{\textgreater{}}\DataTypeTok{m} \DataTypeTok{set\_background\_opacity} \ErrorTok{+}\DataTypeTok{0.1}
\DataTypeTok{map} \DataTypeTok{kitty\_mod}\ErrorTok{+}\DataTypeTok{a}\ErrorTok{\textgreater{}}\DataTypeTok{l} \DataTypeTok{set\_background\_opacity} \DataTypeTok{{-}0.1}
\DataTypeTok{map} \DataTypeTok{kitty\_mod}\ErrorTok{+}\DataTypeTok{a}\ErrorTok{\textgreater{}}\DataTypeTok{1} \DataTypeTok{set\_background\_opacity} \DataTypeTok{1}
\DataTypeTok{map} \DataTypeTok{kitty\_mod}\ErrorTok{+}\DataTypeTok{a}\ErrorTok{\textgreater{}}\DataTypeTok{d} \DataTypeTok{set\_background\_opacity} \DataTypeTok{default}

\CommentTok{\# Reset terminal}
\DataTypeTok{map} \DataTypeTok{kitty\_mod}\ErrorTok{+}\DataTypeTok{delete} \DataTypeTok{clear\_terminal} \DataTypeTok{reset} \DataTypeTok{active}
\end{Highlighting}
\end{Shaded}

\subsection{Configuración Avanzada}\label{configuraciuxf3n-avanzada}

\textbf{Fuentes con fallbacks:}

\begin{Shaded}
\begin{Highlighting}[]
\DataTypeTok{font\_family}      \DataTypeTok{JetBrains} \DataTypeTok{Mono}
\CommentTok{\# Símbolos especiales}
\DataTypeTok{symbol\_map} \DataTypeTok{U}\ErrorTok{+}\DataTypeTok{E0A0{-}U}\ErrorTok{+}\DataTypeTok{E0A3}\ErrorTok{,}\DataTypeTok{U}\ErrorTok{+}\DataTypeTok{E0C0{-}U}\ErrorTok{+}\DataTypeTok{E0C7} \DataTypeTok{PowerlineSymbols}
\CommentTok{\# Emojis}
\DataTypeTok{symbol\_map} \DataTypeTok{U}\ErrorTok{+}\DataTypeTok{1F300{-}U}\ErrorTok{+}\DataTypeTok{1F9FF} \DataTypeTok{Noto} \DataTypeTok{Color} \DataTypeTok{Emoji}
\end{Highlighting}
\end{Shaded}

\textbf{Ligaduras:}

\begin{Shaded}
\begin{Highlighting}[]
\CommentTok{\# Habilitar ligaduras}
\DataTypeTok{disable\_ligatures} \DataTypeTok{never}
\CommentTok{\# O solo cuando el cursor no está sobre ellas}
\CommentTok{\# disable\_ligatures cursor}
\end{Highlighting}
\end{Shaded}

\textbf{Modificar fuentes:}

\begin{Shaded}
\begin{Highlighting}[]
\CommentTok{\# Ajustar posición del underline}
\DataTypeTok{modify\_font} \DataTypeTok{underline\_position} \DataTypeTok{2}
\DataTypeTok{modify\_font} \DataTypeTok{underline\_thickness} \DataTypeTok{150}\ErrorTok{\%}
\CommentTok{\# Ajustar tamaño de celda}
\DataTypeTok{modify\_font} \DataTypeTok{cell\_width} \DataTypeTok{100}\ErrorTok{\%}
\DataTypeTok{modify\_font} \DataTypeTok{cell\_height} \DataTypeTok{100}\ErrorTok{\%}
\end{Highlighting}
\end{Shaded}

\section{Atajos de Teclado}\label{atajos-de-teclado}

\subsection{Modificador Principal}\label{modificador-principal}

Por defecto \texttt{kitty\_mod} = \texttt{Ctrl+Shift}

En macOS algunos también usan \texttt{Cmd} (\texttt{⌘})

\subsection{Scrolling}\label{scrolling}

\begin{longtable}[]{@{}ll@{}}
\toprule\noalign{}
Atajo & Acción \\
\midrule\noalign{}
\endhead
\bottomrule\noalign{}
\endlastfoot
\texttt{Ctrl+Shift+Up} & Scroll línea arriba \\
\texttt{Ctrl+Shift+Down} & Scroll línea abajo \\
\texttt{Ctrl+Shift+Page\ Up} & Scroll página arriba \\
\texttt{Ctrl+Shift+Page\ Down} & Scroll página abajo \\
\texttt{Ctrl+Shift+Home} & Ir al inicio \\
\texttt{Ctrl+Shift+End} & Ir al final \\
\texttt{Ctrl+Shift+Z} & Prompt anterior (shell integration) \\
\texttt{Ctrl+Shift+X} & Prompt siguiente (shell integration) \\
\texttt{Ctrl+Shift+H} & Ver scrollback en pager \\
\texttt{Ctrl+Shift+G} & Ver último comando (shell integration) \\
\texttt{Ctrl+Shift+/} & Buscar en scrollback \\
\end{longtable}

\subsection{Tabs}\label{tabs}

\begin{longtable}[]{@{}ll@{}}
\toprule\noalign{}
Atajo & Acción \\
\midrule\noalign{}
\endhead
\bottomrule\noalign{}
\endlastfoot
\texttt{Ctrl+Shift+T} & Nueva tab \\
\texttt{Ctrl+Shift+Q} & Cerrar tab \\
\texttt{Ctrl+Shift+Right} & Tab siguiente \\
\texttt{Ctrl+Shift+Left} & Tab anterior \\
\texttt{Ctrl+Shift+.} & Mover tab adelante \\
\texttt{Ctrl+Shift+,} & Mover tab atrás \\
\texttt{Ctrl+Shift+Alt+T} & Renombrar tab \\
\texttt{Ctrl+Alt+1...9} & Ir a tab N \\
\end{longtable}

\subsection{Windows (Ventanas dentro de
Tabs)}\label{windows-ventanas-dentro-de-tabs}

\begin{longtable}[]{@{}
  >{\raggedright\arraybackslash}p{(\linewidth - 2\tabcolsep) * \real{0.4667}}
  >{\raggedright\arraybackslash}p{(\linewidth - 2\tabcolsep) * \real{0.5333}}@{}}
\toprule\noalign{}
\begin{minipage}[b]{\linewidth}\raggedright
Atajo
\end{minipage} & \begin{minipage}[b]{\linewidth}\raggedright
Acción
\end{minipage} \\
\midrule\noalign{}
\endhead
\bottomrule\noalign{}
\endlastfoot
\texttt{Ctrl+Shift+Enter} & Nueva window \\
\texttt{Ctrl+Shift+N} & Nueva OS window \\
\texttt{Ctrl+Shift+W} & Cerrar window \\
\texttt{Ctrl+Shift+{]}} & Window siguiente \\
\texttt{Ctrl+Shift+{[}} & Window anterior \\
\texttt{Ctrl+Shift+F} & Mover window adelante \\
\texttt{Ctrl+Shift+B} & Mover window atrás \\
\texttt{Ctrl+Shift+\textasciigrave{}\textasciigrave{}}
`\texttt{\textbar{}\ Mover\ window\ al\ top\ \textbar{}\ \textbar{}}Ctrl+Shift+R\texttt{\textbar{}\ Redimensionar\ window\ \textbar{}\ \textbar{}}Ctrl+Shift+F7\texttt{\textbar{}\ Selección\ visual\ de\ window\ \textbar{}\ \textbar{}}Ctrl+Shift+F8\texttt{\textbar{}\ Swap\ visual\ de\ window\ \textbar{}\ \textbar{}}Ctrl+Shift+1\ldots0`
& Ir a window N \\
\end{longtable}

\subsection{Navegación Direccional}\label{navegaciuxf3n-direccional}

\textbf{Configurar manualmente:}

\begin{Shaded}
\begin{Highlighting}[]
\CommentTok{\# Navegar entre windows}
\DataTypeTok{map} \DataTypeTok{ctrl}\ErrorTok{+}\DataTypeTok{left} \DataTypeTok{neighboring\_window} \DataTypeTok{left}
\DataTypeTok{map} \DataTypeTok{ctrl}\ErrorTok{+}\DataTypeTok{right} \DataTypeTok{neighboring\_window} \DataTypeTok{right}
\DataTypeTok{map} \DataTypeTok{ctrl}\ErrorTok{+}\DataTypeTok{up} \DataTypeTok{neighboring\_window} \DataTypeTok{up}
\DataTypeTok{map} \DataTypeTok{ctrl}\ErrorTok{+}\DataTypeTok{down} \DataTypeTok{neighboring\_window} \DataTypeTok{down}

\CommentTok{\# Mover windows}
\DataTypeTok{map} \DataTypeTok{shift}\ErrorTok{+}\DataTypeTok{left} \DataTypeTok{move\_window} \DataTypeTok{right}
\DataTypeTok{map} \DataTypeTok{shift}\ErrorTok{+}\DataTypeTok{right} \DataTypeTok{move\_window} \DataTypeTok{left}
\DataTypeTok{map} \DataTypeTok{shift}\ErrorTok{+}\DataTypeTok{up} \DataTypeTok{move\_window} \DataTypeTok{down}
\DataTypeTok{map} \DataTypeTok{shift}\ErrorTok{+}\DataTypeTok{down} \DataTypeTok{move\_window} \DataTypeTok{up}
\end{Highlighting}
\end{Shaded}

\subsection{Layouts}\label{layouts}

\begin{longtable}[]{@{}ll@{}}
\toprule\noalign{}
Atajo & Acción \\
\midrule\noalign{}
\endhead
\bottomrule\noalign{}
\endlastfoot
\texttt{Ctrl+Shift+L} & Siguiente layout \\
\end{longtable}

\subsection{Clipboard}\label{clipboard}

\begin{longtable}[]{@{}ll@{}}
\toprule\noalign{}
Atajo & Acción \\
\midrule\noalign{}
\endhead
\bottomrule\noalign{}
\endlastfoot
\texttt{Ctrl+Shift+C} & Copiar \\
\texttt{Ctrl+Shift+V} & Pegar \\
\texttt{Ctrl+Shift+S} & Pegar de selección \\
\texttt{Ctrl+Shift+O} & Pasar selección a programa \\
\end{longtable}

\subsection{Fuentes}\label{fuentes}

\begin{longtable}[]{@{}ll@{}}
\toprule\noalign{}
Atajo & Acción \\
\midrule\noalign{}
\endhead
\bottomrule\noalign{}
\endlastfoot
\texttt{Ctrl+Shift+Equal} & Aumentar tamaño \\
\texttt{Ctrl+Shift+Minus} & Disminuir tamaño \\
\texttt{Ctrl+Shift+Backspace} & Reset tamaño \\
\end{longtable}

\subsection{Miscelánea}\label{misceluxe1nea}

\begin{longtable}[]{@{}ll@{}}
\toprule\noalign{}
Atajo & Acción \\
\midrule\noalign{}
\endhead
\bottomrule\noalign{}
\endlastfoot
\texttt{Ctrl+Shift+F1} & Mostrar ayuda \\
\texttt{Ctrl+Shift+F2} & Editar config \\
\texttt{Ctrl+Shift+F5} & Recargar config \\
\texttt{Ctrl+Shift+F6} & Debug config \\
\texttt{Ctrl+Shift+F11} & Toggle fullscreen \\
\texttt{Ctrl+Shift+F10} & Toggle maximized \\
\texttt{Ctrl+Shift+U} & Input Unicode \\
\texttt{Ctrl+Shift+E} & Open URL hints \\
\texttt{Ctrl+Shift+Escape} & Kitty shell \\
\texttt{Ctrl+Shift+Delete} & Reset terminal \\
\end{longtable}

\subsection{Atajos en macOS}\label{atajos-en-macos}

Reemplazan algunos con \texttt{Cmd} (⌘):

\begin{longtable}[]{@{}ll@{}}
\toprule\noalign{}
Windows/Linux & macOS \\
\midrule\noalign{}
\endhead
\bottomrule\noalign{}
\endlastfoot
\texttt{Ctrl+Shift+T} & \texttt{Cmd+T} \\
\texttt{Ctrl+Shift+W} & \texttt{Cmd+W} \\
\texttt{Ctrl+Shift+C} & \texttt{Cmd+C} \\
\texttt{Ctrl+Shift+V} & \texttt{Cmd+V} \\
\texttt{Ctrl+Shift+N} & \texttt{Cmd+N} \\
\texttt{Ctrl+Shift+F2} & \texttt{Cmd+,} \\
\texttt{Ctrl+Shift+F11} & \texttt{Ctrl+Cmd+F} \\
\end{longtable}

\section{Ventanas y Pestañas}\label{ventanas-y-pestauxf1as}

\subsection{Concepto de
Multiplexación}\label{concepto-de-multiplexaciuxf3n}

Kitty tiene multiplexación nativa (sin necesidad de tmux):

\textbf{Jerarquía:}

\begin{verbatim}
OS Window (kitty instance)
  └── Tab 1
      ├── Window 1
      ├── Window 2
      └── Window 3
  └── Tab 2
      └── Window 1
\end{verbatim}

\subsection{Crear Ventanas y Tabs}\label{crear-ventanas-y-tabs}

\textbf{Nueva window:}

\begin{Shaded}
\begin{Highlighting}[]
\CommentTok{\# Atajo}
\ExtensionTok{Ctrl+Shift+Enter}

\CommentTok{\# Comando}
\ExtensionTok{kitty}\NormalTok{ @ launch }\AttributeTok{{-}{-}type}\OperatorTok{=}\NormalTok{window}

\CommentTok{\# Con directorio específico}
\ExtensionTok{map}\NormalTok{ f1 launch }\AttributeTok{{-}{-}cwd}\OperatorTok{=}\NormalTok{current}
\end{Highlighting}
\end{Shaded}

\textbf{Nueva tab:}

\begin{Shaded}
\begin{Highlighting}[]
\CommentTok{\# Atajo}
\ExtensionTok{Ctrl+Shift+T}

\CommentTok{\# Con programa}
\ExtensionTok{map}\NormalTok{ f1 new\_tab\_with\_cwd}
\end{Highlighting}
\end{Shaded}

\textbf{Nueva OS window:}

\begin{Shaded}
\begin{Highlighting}[]
\ExtensionTok{Ctrl+Shift+N}
\end{Highlighting}
\end{Shaded}

\subsection{Navegación}\label{navegaciuxf3n}

\textbf{Entre windows:}

\begin{Shaded}
\begin{Highlighting}[]
\CommentTok{\# Direccional}
\DataTypeTok{map} \DataTypeTok{ctrl}\ErrorTok{+}\DataTypeTok{left} \DataTypeTok{neighboring\_window} \DataTypeTok{left}
\DataTypeTok{map} \DataTypeTok{ctrl}\ErrorTok{+}\DataTypeTok{right} \DataTypeTok{neighboring\_window} \DataTypeTok{right}
\DataTypeTok{map} \DataTypeTok{ctrl}\ErrorTok{+}\DataTypeTok{up} \DataTypeTok{neighboring\_window} \DataTypeTok{up}
\DataTypeTok{map} \DataTypeTok{ctrl}\ErrorTok{+}\DataTypeTok{down} \DataTypeTok{neighboring\_window} \DataTypeTok{down}

\CommentTok{\# Numérica}
\DataTypeTok{map} \DataTypeTok{ctrl}\ErrorTok{+}\DataTypeTok{shift}\ErrorTok{+}\DataTypeTok{1} \DataTypeTok{first\_window}
\DataTypeTok{map} \DataTypeTok{ctrl}\ErrorTok{+}\DataTypeTok{shift}\ErrorTok{+}\DataTypeTok{2} \DataTypeTok{second\_window}
\CommentTok{\# ...}
\end{Highlighting}
\end{Shaded}

\textbf{Entre tabs:}

\begin{Shaded}
\begin{Highlighting}[]
\DataTypeTok{map} \DataTypeTok{ctrl}\ErrorTok{+}\DataTypeTok{shift}\ErrorTok{+}\DataTypeTok{right} \DataTypeTok{next\_tab}
\DataTypeTok{map} \DataTypeTok{ctrl}\ErrorTok{+}\DataTypeTok{shift}\ErrorTok{+}\DataTypeTok{left} \DataTypeTok{previous\_tab}
\DataTypeTok{map} \DataTypeTok{ctrl}\ErrorTok{+}\DataTypeTok{alt}\ErrorTok{+}\DataTypeTok{1} \DataTypeTok{goto\_tab} \DataTypeTok{1}
\end{Highlighting}
\end{Shaded}

\subsection{Mover Windows y Tabs}\label{mover-windows-y-tabs}

\textbf{Mover windows:}

\begin{Shaded}
\begin{Highlighting}[]
\DataTypeTok{map} \DataTypeTok{ctrl}\ErrorTok{+}\DataTypeTok{shift}\ErrorTok{+}\DataTypeTok{f} \DataTypeTok{move\_window\_forward}
\DataTypeTok{map} \DataTypeTok{ctrl}\ErrorTok{+}\DataTypeTok{shift}\ErrorTok{+}\DataTypeTok{b} \DataTypeTok{move\_window\_backward}
\DataTypeTok{map} \DataTypeTok{ctrl}\ErrorTok{+}\DataTypeTok{shift}\ErrorTok{+\textasciigrave{}} \DataTypeTok{move\_window\_to\_top}

\CommentTok{\# Direccional}
\DataTypeTok{map} \DataTypeTok{shift}\ErrorTok{+}\DataTypeTok{up} \DataTypeTok{move\_window} \DataTypeTok{up}
\DataTypeTok{map} \DataTypeTok{shift}\ErrorTok{+}\DataTypeTok{down} \DataTypeTok{move\_window} \DataTypeTok{down}
\DataTypeTok{map} \DataTypeTok{shift}\ErrorTok{+}\DataTypeTok{left} \DataTypeTok{move\_window} \DataTypeTok{left}
\DataTypeTok{map} \DataTypeTok{shift}\ErrorTok{+}\DataTypeTok{right} \DataTypeTok{move\_window} \DataTypeTok{right}
\end{Highlighting}
\end{Shaded}

\textbf{Mover tabs:}

\begin{Shaded}
\begin{Highlighting}[]
\DataTypeTok{map} \DataTypeTok{ctrl}\ErrorTok{+}\DataTypeTok{shift}\ErrorTok{+}\DataTypeTok{.} \DataTypeTok{move\_tab\_forward}
\DataTypeTok{map} \DataTypeTok{ctrl}\ErrorTok{+}\DataTypeTok{shift}\ErrorTok{+,} \DataTypeTok{move\_tab\_backward}
\end{Highlighting}
\end{Shaded}

\subsection{Detach (Separar)}\label{detach-separar}

\textbf{Detach window:}

\begin{Shaded}
\begin{Highlighting}[]
\CommentTok{\# A nueva OS window}
\DataTypeTok{map} \DataTypeTok{ctrl}\ErrorTok{+}\DataTypeTok{f2} \DataTypeTok{detach\_window}

\CommentTok{\# A nueva tab}
\DataTypeTok{map} \DataTypeTok{ctrl}\ErrorTok{+}\DataTypeTok{f3} \DataTypeTok{detach\_window} \DataTypeTok{new{-}tab}

\CommentTok{\# A tab anterior}
\DataTypeTok{map} \DataTypeTok{ctrl}\ErrorTok{+}\DataTypeTok{f3} \DataTypeTok{detach\_window} \DataTypeTok{tab{-}prev}

\CommentTok{\# Preguntar destino}
\DataTypeTok{map} \DataTypeTok{ctrl}\ErrorTok{+}\DataTypeTok{f4} \DataTypeTok{detach\_window} \DataTypeTok{ask}
\end{Highlighting}
\end{Shaded}

\textbf{Detach tab:}

\begin{Shaded}
\begin{Highlighting}[]
\CommentTok{\# A nueva OS window}
\DataTypeTok{map} \DataTypeTok{ctrl}\ErrorTok{+}\DataTypeTok{f2} \DataTypeTok{detach\_tab}

\CommentTok{\# Preguntar destino}
\DataTypeTok{map} \DataTypeTok{ctrl}\ErrorTok{+}\DataTypeTok{f4} \DataTypeTok{detach\_tab} \DataTypeTok{ask}
\end{Highlighting}
\end{Shaded}

\subsection{Títulos Personalizados}\label{tuxedtulos-personalizados}

\begin{Shaded}
\begin{Highlighting}[]
\CommentTok{\# Renombrar tab}
\DataTypeTok{map} \DataTypeTok{ctrl}\ErrorTok{+}\DataTypeTok{shift}\ErrorTok{+}\DataTypeTok{alt}\ErrorTok{+}\DataTypeTok{t} \DataTypeTok{set\_tab\_title}

\CommentTok{\# Renombrar window}
\DataTypeTok{map} \DataTypeTok{f1} \DataTypeTok{set\_window\_title}

\CommentTok{\# Template de título}
\DataTypeTok{tab\_title\_template} \DataTypeTok{"\{index\}: \{title\}"}
\end{Highlighting}
\end{Shaded}

\section{Layouts (Diseños)}\label{layouts-diseuxf1os}

\subsection{Layouts Disponibles}\label{layouts-disponibles}

Kitty tiene 7 layouts built-in:

\begin{enumerate}
\def\labelenumi{\arabic{enumi}.}
\tightlist
\item
  \textbf{Stack}: Una window a la vez (maximizada)
\item
  \textbf{Tall}: Una o más windows full-height a la izquierda
\item
  \textbf{Fat}: Una o más windows full-width arriba
\item
  \textbf{Grid}: Grid balanceado
\item
  \textbf{Splits}: Arreglo arbitrario (como tmux)
\item
  \textbf{Horizontal}: Todas las windows lado a lado
\item
  \textbf{Vertical}: Todas las windows una debajo de otra
\end{enumerate}

\subsection{Stack Layout}\label{stack-layout}

\begin{verbatim}
┌────────────────────────────┐
│                            │
│                            │
│       Window Actual        │
│       (maximizada)         │
│                            │
│                            │
└────────────────────────────┘
\end{verbatim}

\textbf{Configuración:}

\begin{Shaded}
\begin{Highlighting}[]
\DataTypeTok{enabled\_layouts} \DataTypeTok{stack}
\end{Highlighting}
\end{Shaded}

\subsection{Tall Layout}\label{tall-layout}

\begin{verbatim}
┌──────────┬─────────────────┐
│          │                 │
│          │                 │
│          ├─────────────────┤
│  Full    │                 │
│  Height  │                 │
│          ├─────────────────┤
│          │                 │
│          │                 │
└──────────┴─────────────────┘
\end{verbatim}

\textbf{Configuración:}

\begin{Shaded}
\begin{Highlighting}[]
\DataTypeTok{enabled\_layouts} \DataTypeTok{tall}\ErrorTok{:}\DataTypeTok{bias}\OperatorTok{=}\DecValTok{50}\ErrorTok{;}\DataTypeTok{full\_size}\OperatorTok{=}\DecValTok{1}\ErrorTok{;}\DataTypeTok{mirrored}\OperatorTok{=}\ConstantTok{false}

\CommentTok{\# Opciones:}
\CommentTok{\# bias: 10{-}90 (split horizontal)}
\CommentTok{\# full\_size: número de windows full{-}height}
\CommentTok{\# mirrored: true/false (full{-}height a la derecha)}
\end{Highlighting}
\end{Shaded}

\textbf{Acciones:}

\begin{Shaded}
\begin{Highlighting}[]
\CommentTok{\# Aumentar/disminuir windows full{-}height}
\DataTypeTok{map} \DataTypeTok{ctrl}\ErrorTok{+}\KeywordTok{[} \KeywordTok{layout\_action} \KeywordTok{decrease\_num\_full\_size\_windows}
\DataTypeTok{map} \DataTypeTok{ctrl}\ErrorTok{+]} \DataTypeTok{layout\_action} \DataTypeTok{increase\_num\_full\_size\_windows}

\CommentTok{\# Toggle mirror}
\DataTypeTok{map} \DataTypeTok{ctrl}\ErrorTok{+/} \DataTypeTok{layout\_action} \DataTypeTok{mirror} \DataTypeTok{toggle}

\CommentTok{\# Cambiar bias}
\DataTypeTok{map} \DataTypeTok{ctrl}\ErrorTok{+}\DataTypeTok{.} \DataTypeTok{layout\_action} \DataTypeTok{bias} \DataTypeTok{50} \DataTypeTok{62} \DataTypeTok{70}
\end{Highlighting}
\end{Shaded}

\subsection{Fat Layout}\label{fat-layout}

\begin{verbatim}
┌────────────────────────────┐
│                            │
│      Full Width            │
│                            │
├─────────┬─────────┬────────┤
│         │         │        │
│         │         │        │
└─────────┴─────────┴────────┘
\end{verbatim}

\textbf{Configuración:}

\begin{Shaded}
\begin{Highlighting}[]
\DataTypeTok{enabled\_layouts} \DataTypeTok{fat}\ErrorTok{:}\DataTypeTok{bias}\OperatorTok{=}\DecValTok{50}\ErrorTok{;}\DataTypeTok{full\_size}\OperatorTok{=}\DecValTok{1}\ErrorTok{;}\DataTypeTok{mirrored}\OperatorTok{=}\ConstantTok{false}
\end{Highlighting}
\end{Shaded}

Mismas acciones que Tall layout.

\subsection{Grid Layout}\label{grid-layout}

\begin{verbatim}
┌────────┬────────┬────────┐
│        │        │        │
│        │        │        │
├────────┼────────┼────────┤
│        │        │        │
│        │        │        │
└────────┴────────┴────────┘
\end{verbatim}

\textbf{Configuración:}

\begin{Shaded}
\begin{Highlighting}[]
\DataTypeTok{enabled\_layouts} \DataTypeTok{grid}
\end{Highlighting}
\end{Shaded}

\subsection{Splits Layout}\label{splits-layout}

El más flexible - como tmux/vim splits.

\begin{verbatim}
┌──────────┬─────────────────┐
│          │                 │
│          │                 │
│          ├────────┬────────┤
│          │        │        │
│          │        │        │
│          ├────────┴────────┤
│          │                 │
│          │                 │
└──────────┴─────────────────┘
\end{verbatim}

\textbf{Configuración:}

\begin{Shaded}
\begin{Highlighting}[]
\DataTypeTok{enabled\_layouts} \DataTypeTok{splits}\ErrorTok{:}\DataTypeTok{split\_axis}\OperatorTok{=}\DataTypeTok{horizontal}
\CommentTok{\# split\_axis: horizontal, vertical, auto}
\end{Highlighting}
\end{Shaded}

\textbf{Crear splits:}

\begin{Shaded}
\begin{Highlighting}[]
\CommentTok{\# Split horizontal}
\DataTypeTok{map} \DataTypeTok{f5} \DataTypeTok{launch} \DataTypeTok{{-}{-}location}\OperatorTok{=}\DataTypeTok{hsplit}

\CommentTok{\# Split vertical}
\DataTypeTok{map} \DataTypeTok{f6} \DataTypeTok{launch} \DataTypeTok{{-}{-}location}\OperatorTok{=}\DataTypeTok{vsplit}

\CommentTok{\# Split auto (horizontal si ancho, vertical si alto)}
\DataTypeTok{map} \DataTypeTok{f4} \DataTypeTok{launch} \DataTypeTok{{-}{-}location}\OperatorTok{=}\DataTypeTok{split}

\CommentTok{\# Rotar split (horizontal \textless{}{-}\textgreater{} vertical)}
\DataTypeTok{map} \DataTypeTok{f7} \DataTypeTok{layout\_action} \DataTypeTok{rotate}

\CommentTok{\# Mover window}
\DataTypeTok{map} \DataTypeTok{shift}\ErrorTok{+}\DataTypeTok{up} \DataTypeTok{move\_window} \DataTypeTok{up}
\DataTypeTok{map} \DataTypeTok{shift}\ErrorTok{+}\DataTypeTok{down} \DataTypeTok{move\_window} \DataTypeTok{down}
\DataTypeTok{map} \DataTypeTok{shift}\ErrorTok{+}\DataTypeTok{left} \DataTypeTok{move\_window} \DataTypeTok{left}
\DataTypeTok{map} \DataTypeTok{shift}\ErrorTok{+}\DataTypeTok{right} \DataTypeTok{move\_window} \DataTypeTok{right}

\CommentTok{\# Mover window al borde}
\DataTypeTok{map} \DataTypeTok{ctrl}\ErrorTok{+}\DataTypeTok{shift}\ErrorTok{+}\DataTypeTok{up} \DataTypeTok{layout\_action} \DataTypeTok{move\_to\_screen\_edge} \DataTypeTok{top}
\DataTypeTok{map} \DataTypeTok{ctrl}\ErrorTok{+}\DataTypeTok{shift}\ErrorTok{+}\DataTypeTok{down} \DataTypeTok{layout\_action} \DataTypeTok{move\_to\_screen\_edge} \DataTypeTok{bottom}
\DataTypeTok{map} \DataTypeTok{ctrl}\ErrorTok{+}\DataTypeTok{shift}\ErrorTok{+}\DataTypeTok{left} \DataTypeTok{layout\_action} \DataTypeTok{move\_to\_screen\_edge} \DataTypeTok{left}
\DataTypeTok{map} \DataTypeTok{ctrl}\ErrorTok{+}\DataTypeTok{shift}\ErrorTok{+}\DataTypeTok{right} \DataTypeTok{layout\_action} \DataTypeTok{move\_to\_screen\_edge} \DataTypeTok{right}

\CommentTok{\# Bias del split}
\DataTypeTok{map} \DataTypeTok{ctrl}\ErrorTok{+}\DataTypeTok{.} \DataTypeTok{layout\_action} \DataTypeTok{bias} \DataTypeTok{80}
\end{Highlighting}
\end{Shaded}

\subsection{Horizontal y Vertical
Layouts}\label{horizontal-y-vertical-layouts}

\textbf{Horizontal (lado a lado):}

\begin{verbatim}
┌────────┬────────┬────────┐
│        │        │        │
│        │        │        │
│        │        │        │
│        │        │        │
└────────┴────────┴────────┘
\end{verbatim}

\textbf{Vertical (uno debajo del otro):}

\begin{verbatim}
┌────────────────────────────┐
│                            │
├────────────────────────────┤
│                            │
├────────────────────────────┤
│                            │
└────────────────────────────┘
\end{verbatim}

\textbf{Configuración:}

\begin{Shaded}
\begin{Highlighting}[]
\DataTypeTok{enabled\_layouts} \DataTypeTok{horizontal}\ErrorTok{,}\DataTypeTok{vertical}
\end{Highlighting}
\end{Shaded}

\subsection{Cambiar Entre Layouts}\label{cambiar-entre-layouts}

\begin{Shaded}
\begin{Highlighting}[]
\CommentTok{\# Siguiente layout}
\DataTypeTok{map} \DataTypeTok{ctrl}\ErrorTok{+}\DataTypeTok{shift}\ErrorTok{+}\DataTypeTok{l} \DataTypeTok{next\_layout}

\CommentTok{\# Ir a layout específico}
\DataTypeTok{map} \DataTypeTok{ctrl}\ErrorTok{+}\DataTypeTok{alt}\ErrorTok{+}\DataTypeTok{t} \DataTypeTok{goto\_layout} \DataTypeTok{tall}
\DataTypeTok{map} \DataTypeTok{ctrl}\ErrorTok{+}\DataTypeTok{alt}\ErrorTok{+}\DataTypeTok{s} \DataTypeTok{goto\_layout} \DataTypeTok{stack}

\CommentTok{\# Toggle layout}
\DataTypeTok{map} \DataTypeTok{ctrl}\ErrorTok{+}\DataTypeTok{alt}\ErrorTok{+}\DataTypeTok{z} \DataTypeTok{toggle\_layout} \DataTypeTok{stack}
\end{Highlighting}
\end{Shaded}

\subsection{Redimensionar Windows}\label{redimensionar-windows}

\textbf{Modo interactivo:}

\begin{Shaded}
\begin{Highlighting}[]
\ExtensionTok{Ctrl+Shift+R}
\CommentTok{\# Seguir instrucciones en pantalla}
\end{Highlighting}
\end{Shaded}

\textbf{Atajos directos:}

\begin{Shaded}
\begin{Highlighting}[]
\DataTypeTok{map} \DataTypeTok{ctrl}\ErrorTok{+}\DataTypeTok{left} \DataTypeTok{resize\_window} \DataTypeTok{narrower}
\DataTypeTok{map} \DataTypeTok{ctrl}\ErrorTok{+}\DataTypeTok{right} \DataTypeTok{resize\_window} \DataTypeTok{wider}
\DataTypeTok{map} \DataTypeTok{ctrl}\ErrorTok{+}\DataTypeTok{up} \DataTypeTok{resize\_window} \DataTypeTok{taller}
\DataTypeTok{map} \DataTypeTok{ctrl}\ErrorTok{+}\DataTypeTok{down} \DataTypeTok{resize\_window} \DataTypeTok{shorter} \DataTypeTok{3}
\DataTypeTok{map} \DataTypeTok{ctrl}\ErrorTok{+}\DataTypeTok{home} \DataTypeTok{resize\_window} \DataTypeTok{reset}
\end{Highlighting}
\end{Shaded}

\section{Fuentes}\label{fuentes-1}

\subsection{Configuración Básica}\label{configuraciuxf3n-buxe1sica}

\begin{Shaded}
\begin{Highlighting}[]
\DataTypeTok{font\_family}      \DataTypeTok{JetBrains} \DataTypeTok{Mono}
\DataTypeTok{bold\_font}        \DataTypeTok{auto}
\DataTypeTok{italic\_font}      \DataTypeTok{auto}
\DataTypeTok{bold\_italic\_font} \DataTypeTok{auto}
\DataTypeTok{font\_size}        \DataTypeTok{12.0}
\end{Highlighting}
\end{Shaded}

\subsection{Fuentes Recomendadas}\label{fuentes-recomendadas}

\textbf{Con ligaduras:}

\begin{itemize}
\tightlist
\item
  JetBrains Mono
\item
  Fira Code
\item
  Cascadia Code
\item
  Victor Mono
\item
  Iosevka
\end{itemize}

\textbf{Mono clásicas:}

\begin{itemize}
\tightlist
\item
  IBM Plex Mono
\item
  Source Code Pro
\item
  Ubuntu Mono
\item
  Inconsolata
\end{itemize}

\subsection{Fuentes con Powerline}\label{fuentes-con-powerline}

\begin{Shaded}
\begin{Highlighting}[]
\CommentTok{\# Método 1: Usar font con powerline incluido}
\DataTypeTok{font\_family} \DataTypeTok{JetBrains} \DataTypeTok{Mono} \DataTypeTok{Nerd} \DataTypeTok{Font}

\CommentTok{\# Método 2: Symbol map}
\DataTypeTok{font\_family} \DataTypeTok{JetBrains} \DataTypeTok{Mono}
\DataTypeTok{symbol\_map} \DataTypeTok{U}\ErrorTok{+}\DataTypeTok{E0A0{-}U}\ErrorTok{+}\DataTypeTok{E0A3}\ErrorTok{,}\DataTypeTok{U}\ErrorTok{+}\DataTypeTok{E0C0{-}U}\ErrorTok{+}\DataTypeTok{E0C7} \DataTypeTok{PowerlineSymbols}
\end{Highlighting}
\end{Shaded}

\subsection{Symbol Maps}\label{symbol-maps}

Mapear rangos Unicode a fuentes específicas:

\begin{Shaded}
\begin{Highlighting}[]
\CommentTok{\# Powerline}
\DataTypeTok{symbol\_map} \DataTypeTok{U}\ErrorTok{+}\DataTypeTok{E0A0{-}U}\ErrorTok{+}\DataTypeTok{E0A3}\ErrorTok{,}\DataTypeTok{U}\ErrorTok{+}\DataTypeTok{E0C0{-}U}\ErrorTok{+}\DataTypeTok{E0C7} \DataTypeTok{PowerlineSymbols}

\CommentTok{\# Emojis}
\DataTypeTok{symbol\_map} \DataTypeTok{U}\ErrorTok{+}\DataTypeTok{1F300{-}U}\ErrorTok{+}\DataTypeTok{1F64F}\ErrorTok{,}\DataTypeTok{U}\ErrorTok{+}\DataTypeTok{1F680{-}U}\ErrorTok{+}\DataTypeTok{1F6FF} \DataTypeTok{Noto} \DataTypeTok{Color} \DataTypeTok{Emoji}

\CommentTok{\# Math}
\DataTypeTok{symbol\_map} \DataTypeTok{U}\ErrorTok{+}\DataTypeTok{2200{-}U}\ErrorTok{+}\DataTypeTok{22FF} \DataTypeTok{STIX} \DataTypeTok{Two} \DataTypeTok{Math}
\end{Highlighting}
\end{Shaded}

\subsection{Ligaduras}\label{ligaduras}

\begin{Shaded}
\begin{Highlighting}[]
\CommentTok{\# Siempre habilitadas}
\DataTypeTok{disable\_ligatures} \DataTypeTok{never}

\CommentTok{\# Deshabilitadas cuando cursor está sobre ellas}
\DataTypeTok{disable\_ligatures} \DataTypeTok{cursor}

\CommentTok{\# Siempre deshabilitadas}
\DataTypeTok{disable\_ligatures} \DataTypeTok{always}
\end{Highlighting}
\end{Shaded}

\textbf{Controlar por window:}

\begin{Shaded}
\begin{Highlighting}[]
\DataTypeTok{map} \DataTypeTok{alt}\ErrorTok{+}\DataTypeTok{1} \DataTypeTok{disable\_ligatures\_in} \DataTypeTok{active} \DataTypeTok{always}
\DataTypeTok{map} \DataTypeTok{alt}\ErrorTok{+}\DataTypeTok{2} \DataTypeTok{disable\_ligatures\_in} \DataTypeTok{all} \DataTypeTok{never}
\DataTypeTok{map} \DataTypeTok{alt}\ErrorTok{+}\DataTypeTok{3} \DataTypeTok{disable\_ligatures\_in} \DataTypeTok{tab} \DataTypeTok{cursor}
\end{Highlighting}
\end{Shaded}

\subsection{Font Features}\label{font-features}

\begin{Shaded}
\begin{Highlighting}[]
\CommentTok{\# Habilitar features específicas}
\CommentTok{\# Ejemplo con Fira Code:}
\DataTypeTok{font\_features} \DataTypeTok{FiraCode{-}Regular} \ErrorTok{+}\DataTypeTok{zero} \ErrorTok{+}\DataTypeTok{onum} \ErrorTok{+}\DataTypeTok{ss01} \ErrorTok{+}\DataTypeTok{ss02}

\CommentTok{\# Deshabilitar ligatures pero mantener calt}
\DataTypeTok{font\_features} \DataTypeTok{FiraCode{-}Regular} \DataTypeTok{{-}liga} \ErrorTok{+}\DataTypeTok{calt}
\end{Highlighting}
\end{Shaded}

\subsection{Modificar Fuentes}\label{modificar-fuentes}

\begin{Shaded}
\begin{Highlighting}[]
\CommentTok{\# Ajustar underline}
\DataTypeTok{modify\_font} \DataTypeTok{underline\_position} \DataTypeTok{{-}2}
\DataTypeTok{modify\_font} \DataTypeTok{underline\_thickness} \DataTypeTok{150}\ErrorTok{\%}

\CommentTok{\# Ajustar strikethrough}
\DataTypeTok{modify\_font} \DataTypeTok{strikethrough\_position} \DataTypeTok{2px}

\CommentTok{\# Ajustar tamaño de celda}
\DataTypeTok{modify\_font} \DataTypeTok{cell\_width} \DataTypeTok{80}\ErrorTok{\%}
\DataTypeTok{modify\_font} \DataTypeTok{cell\_height} \DataTypeTok{{-}2px}

\CommentTok{\# Ajustar baseline}
\DataTypeTok{modify\_font} \DataTypeTok{baseline} \DataTypeTok{3}
\end{Highlighting}
\end{Shaded}

\subsection{Kitten para Elegir
Fuentes}\label{kitten-para-elegir-fuentes}

\begin{Shaded}
\begin{Highlighting}[]
\CommentTok{\# UI interactiva para elegir fuentes}
\ExtensionTok{kitten}\NormalTok{ choose{-}fonts}

\CommentTok{\# Previsualizar todas las fuentes}
\ExtensionTok{kitten}\NormalTok{ choose{-}fonts }\AttributeTok{{-}{-}preview}
\end{Highlighting}
\end{Shaded}

\section{Colores y Temas}\label{colores-y-temas}

\subsection{Estructura de Colores}\label{estructura-de-colores}

\begin{Shaded}
\begin{Highlighting}[]
\CommentTok{\# Colores básicos}
\DataTypeTok{foreground} \CommentTok{\#foreground\_color}
\DataTypeTok{background} \CommentTok{\#background\_color}

\CommentTok{\# Selección}
\DataTypeTok{selection\_foreground} \CommentTok{\#color}
\DataTypeTok{selection\_background} \CommentTok{\#color}

\CommentTok{\# Cursor}
\DataTypeTok{cursor} \CommentTok{\#color}
\DataTypeTok{cursor\_text\_color} \CommentTok{\#color}

\CommentTok{\# Tabla de colores (16 básicos)}
\CommentTok{\# Negro}
\DataTypeTok{color0}  \CommentTok{\#color}
\DataTypeTok{color8}  \CommentTok{\#color}

\CommentTok{\# Rojo}
\DataTypeTok{color1}  \CommentTok{\#color}
\DataTypeTok{color9}  \CommentTok{\#color}

\CommentTok{\# Verde}
\DataTypeTok{color2}  \CommentTok{\#color}
\DataTypeTok{color10} \CommentTok{\#color}

\CommentTok{\# Amarillo}
\DataTypeTok{color3}  \CommentTok{\#color}
\DataTypeTok{color11} \CommentTok{\#color}

\CommentTok{\# Azul}
\DataTypeTok{color4}  \CommentTok{\#color}
\DataTypeTok{color12} \CommentTok{\#color}

\CommentTok{\# Magenta}
\DataTypeTok{color5}  \CommentTok{\#color}
\DataTypeTok{color13} \CommentTok{\#color}

\CommentTok{\# Cyan}
\DataTypeTok{color6}  \CommentTok{\#color}
\DataTypeTok{color14} \CommentTok{\#color}

\CommentTok{\# Blanco}
\DataTypeTok{color7}  \CommentTok{\#color}
\DataTypeTok{color15} \CommentTok{\#color}

\CommentTok{\# Colores extendidos (256 colores)}
\DataTypeTok{color16} \CommentTok{\#color}
\DataTypeTok{color17} \CommentTok{\#color}
\CommentTok{\# ... hasta color255}
\end{Highlighting}
\end{Shaded}

\subsection{Temas Populares}\label{temas-populares}

\textbf{Dracula:}

\begin{Shaded}
\begin{Highlighting}[]
\DataTypeTok{foreground} \CommentTok{\#f8f8f2}
\DataTypeTok{background} \CommentTok{\#282a36}
\DataTypeTok{selection\_foreground} \CommentTok{\#ffffff}
\DataTypeTok{selection\_background} \CommentTok{\#44475a}

\DataTypeTok{color0}  \CommentTok{\#21222c}
\DataTypeTok{color8}  \CommentTok{\#6272a4}
\DataTypeTok{color1}  \CommentTok{\#ff5555}
\DataTypeTok{color9}  \CommentTok{\#ff6e6e}
\DataTypeTok{color2}  \CommentTok{\#50fa7b}
\DataTypeTok{color10} \CommentTok{\#69ff94}
\DataTypeTok{color3}  \CommentTok{\#f1fa8c}
\DataTypeTok{color11} \CommentTok{\#ffffa5}
\DataTypeTok{color4}  \CommentTok{\#bd93f9}
\DataTypeTok{color12} \CommentTok{\#d6acff}
\DataTypeTok{color5}  \CommentTok{\#ff79c6}
\DataTypeTok{color13} \CommentTok{\#ff92df}
\DataTypeTok{color6}  \CommentTok{\#8be9fd}
\DataTypeTok{color14} \CommentTok{\#a4ffff}
\DataTypeTok{color7}  \CommentTok{\#f8f8f2}
\DataTypeTok{color15} \CommentTok{\#ffffff}
\end{Highlighting}
\end{Shaded}

\textbf{Nord:}

\begin{Shaded}
\begin{Highlighting}[]
\DataTypeTok{foreground} \CommentTok{\#d8dee9}
\DataTypeTok{background} \CommentTok{\#2e3440}

\DataTypeTok{color0}  \CommentTok{\#3b4252}
\DataTypeTok{color8}  \CommentTok{\#4c566a}
\DataTypeTok{color1}  \CommentTok{\#bf616a}
\DataTypeTok{color9}  \CommentTok{\#bf616a}
\DataTypeTok{color2}  \CommentTok{\#a3be8c}
\DataTypeTok{color10} \CommentTok{\#a3be8c}
\DataTypeTok{color3}  \CommentTok{\#ebcb8b}
\DataTypeTok{color11} \CommentTok{\#ebcb8b}
\DataTypeTok{color4}  \CommentTok{\#81a1c1}
\DataTypeTok{color12} \CommentTok{\#81a1c1}
\DataTypeTok{color5}  \CommentTok{\#b48ead}
\DataTypeTok{color13} \CommentTok{\#b48ead}
\DataTypeTok{color6}  \CommentTok{\#88c0d0}
\DataTypeTok{color14} \CommentTok{\#8fbcbb}
\DataTypeTok{color7}  \CommentTok{\#e5e9f0}
\DataTypeTok{color15} \CommentTok{\#eceff4}
\end{Highlighting}
\end{Shaded}

\textbf{Gruvbox Dark:}

\begin{Shaded}
\begin{Highlighting}[]
\DataTypeTok{foreground} \CommentTok{\#ebdbb2}
\DataTypeTok{background} \CommentTok{\#282828}

\DataTypeTok{color0}  \CommentTok{\#282828}
\DataTypeTok{color8}  \CommentTok{\#928374}
\DataTypeTok{color1}  \CommentTok{\#cc241d}
\DataTypeTok{color9}  \CommentTok{\#fb4934}
\DataTypeTok{color2}  \CommentTok{\#98971a}
\DataTypeTok{color10} \CommentTok{\#b8bb26}
\DataTypeTok{color3}  \CommentTok{\#d79921}
\DataTypeTok{color11} \CommentTok{\#fabd2f}
\DataTypeTok{color4}  \CommentTok{\#458588}
\DataTypeTok{color12} \CommentTok{\#83a598}
\DataTypeTok{color5}  \CommentTok{\#b16286}
\DataTypeTok{color13} \CommentTok{\#d3869b}
\DataTypeTok{color6}  \CommentTok{\#689d6a}
\DataTypeTok{color14} \CommentTok{\#8ec07c}
\DataTypeTok{color7}  \CommentTok{\#a89984}
\DataTypeTok{color15} \CommentTok{\#ebdbb2}
\end{Highlighting}
\end{Shaded}

\subsection{Kitten Themes}\label{kitten-themes}

\begin{Shaded}
\begin{Highlighting}[]
\CommentTok{\# Browser de temas interactivo}
\ExtensionTok{kitten}\NormalTok{ themes}

\CommentTok{\# Listar temas disponibles}
\ExtensionTok{kitten}\NormalTok{ themes }\AttributeTok{{-}{-}list}

\CommentTok{\# Aplicar tema específico}
\ExtensionTok{kitten}\NormalTok{ themes }\AttributeTok{{-}{-}reload{-}in}\OperatorTok{=}\NormalTok{all Dracula}
\end{Highlighting}
\end{Shaded}

\subsection{Opacidad y Blur}\label{opacidad-y-blur}

\begin{Shaded}
\begin{Highlighting}[]
\CommentTok{\# Opacidad del background}
\DataTypeTok{background\_opacity} \DataTypeTok{0.9}

\CommentTok{\# Background blur (macOS y KDE)}
\DataTypeTok{background\_blur} \DataTypeTok{20}

\CommentTok{\# Dynamic opacity (permite cambiar con atajos)}
\DataTypeTok{dynamic\_background\_opacity} \DataTypeTok{yes}

\CommentTok{\# Cambiar opacidad con atajos}
\DataTypeTok{map} \DataTypeTok{ctrl}\ErrorTok{+}\DataTypeTok{shift}\ErrorTok{+}\DataTypeTok{a}\ErrorTok{\textgreater{}}\DataTypeTok{m} \DataTypeTok{set\_background\_opacity} \ErrorTok{+}\DataTypeTok{0.1}
\DataTypeTok{map} \DataTypeTok{ctrl}\ErrorTok{+}\DataTypeTok{shift}\ErrorTok{+}\DataTypeTok{a}\ErrorTok{\textgreater{}}\DataTypeTok{l} \DataTypeTok{set\_background\_opacity} \DataTypeTok{{-}0.1}
\DataTypeTok{map} \DataTypeTok{ctrl}\ErrorTok{+}\DataTypeTok{shift}\ErrorTok{+}\DataTypeTok{a}\ErrorTok{\textgreater{}}\DataTypeTok{1} \DataTypeTok{set\_background\_opacity} \DataTypeTok{1}
\DataTypeTok{map} \DataTypeTok{ctrl}\ErrorTok{+}\DataTypeTok{shift}\ErrorTok{+}\DataTypeTok{a}\ErrorTok{\textgreater{}}\DataTypeTok{d} \DataTypeTok{set\_background\_opacity} \DataTypeTok{default}
\end{Highlighting}
\end{Shaded}

\subsection{Background Image}\label{background-image}

\begin{Shaded}
\begin{Highlighting}[]
\DataTypeTok{background\_image} \ErrorTok{/}\DataTypeTok{path}\ErrorTok{/}\DataTypeTok{to}\ErrorTok{/}\DataTypeTok{image.png}
\DataTypeTok{background\_image\_layout} \DataTypeTok{tiled}
\CommentTok{\# Opciones: tiled, mirror{-}tiled, scaled, clamped, centered, cscaled}
\DataTypeTok{background\_image\_linear} \DataTypeTok{no}
\DataTypeTok{background\_tint} \DataTypeTok{0.5}
\end{Highlighting}
\end{Shaded}

\subsection{Dim Text}\label{dim-text}

\begin{Shaded}
\begin{Highlighting}[]
\CommentTok{\# Atenuar texto con atributo DIM}
\DataTypeTok{dim\_opacity} \DataTypeTok{0.4}
\end{Highlighting}
\end{Shaded}

\section{Cursor y Scrollback}\label{cursor-y-scrollback}

\subsection{Cursor}\label{cursor}

\textbf{Forma:}

\begin{Shaded}
\begin{Highlighting}[]
\DataTypeTok{cursor\_shape} \DataTypeTok{block}
\CommentTok{\# Opciones: block, beam, underline}
\end{Highlighting}
\end{Shaded}

\textbf{Forma cuando no enfocado:}

\begin{Shaded}
\begin{Highlighting}[]
\DataTypeTok{cursor\_shape\_unfocused} \DataTypeTok{hollow}
\CommentTok{\# Opciones: block, beam, underline, hollow, unchanged}
\end{Highlighting}
\end{Shaded}

\textbf{Grosor:}

\begin{Shaded}
\begin{Highlighting}[]
\DataTypeTok{cursor\_beam\_thickness} \DataTypeTok{1.5}
\DataTypeTok{cursor\_underline\_thickness} \DataTypeTok{2.0}
\end{Highlighting}
\end{Shaded}

\textbf{Blink:}

\begin{Shaded}
\begin{Highlighting}[]
\CommentTok{\# Intervalo de blink (segundos)}
\DataTypeTok{cursor\_blink\_interval} \DataTypeTok{{-}1}
\CommentTok{\# {-}1 = usar default del sistema}
\CommentTok{\# 0 = deshabilitado}
\CommentTok{\# \textgreater{}0 = intervalo personalizado}

\CommentTok{\# Parar blink después de inactividad}
\DataTypeTok{cursor\_stop\_blinking\_after} \DataTypeTok{15.0}

\CommentTok{\# Animación de blink}
\DataTypeTok{cursor\_blink\_interval} \DataTypeTok{0.5} \DataTypeTok{ease{-}in{-}out}
\end{Highlighting}
\end{Shaded}

\textbf{Color:}

\begin{Shaded}
\begin{Highlighting}[]
\DataTypeTok{cursor} \CommentTok{\#cccccc}
\DataTypeTok{cursor\_text\_color} \CommentTok{\#111111}
\CommentTok{\# O usar: cursor\_text\_color background}
\end{Highlighting}
\end{Shaded}

\textbf{Cursor Trail:}

\begin{Shaded}
\begin{Highlighting}[]
\CommentTok{\# Animación de "estela" del cursor}
\DataTypeTok{cursor\_trail} \DataTypeTok{5}
\DataTypeTok{cursor\_trail\_decay} \DataTypeTok{0.1} \DataTypeTok{0.4}
\DataTypeTok{cursor\_trail\_start\_threshold} \DataTypeTok{2}
\DataTypeTok{cursor\_trail\_color} \CommentTok{\#ff0000}
\end{Highlighting}
\end{Shaded}

\subsection{Scrollback}\label{scrollback}

\textbf{Líneas de historial:}

\begin{Shaded}
\begin{Highlighting}[]
\DataTypeTok{scrollback\_lines} \DataTypeTok{10000}
\CommentTok{\# {-}1 = infinito (no recomendado)}
\end{Highlighting}
\end{Shaded}

\textbf{Pager:}

\begin{Shaded}
\begin{Highlighting}[]
\DataTypeTok{scrollback\_pager} \DataTypeTok{less} \DataTypeTok{{-}{-}chop{-}long{-}lines} \DataTypeTok{{-}{-}RAW{-}CONTROL{-}CHARS} \ErrorTok{+}\DataTypeTok{INPUT\_LINE\_NUMBER}

\CommentTok{\# Con neovim}
\DataTypeTok{scrollback\_pager} \DataTypeTok{nvim} \DataTypeTok{{-}{-}cmd} \DataTypeTok{\textquotesingle{}set eventignore=FileType\textquotesingle{}} \ErrorTok{+}\DataTypeTok{\textquotesingle{}nnoremap q ZQ\textquotesingle{}} \ErrorTok{+}\DataTypeTok{\textquotesingle{}call nvim\_open\_term(0, \{\})\textquotesingle{}} \ErrorTok{+}\DataTypeTok{\textquotesingle{}set nomodified nolist\textquotesingle{}} \ErrorTok{+}\DataTypeTok{\textquotesingle{}$\textquotesingle{}} \DataTypeTok{{-}}
\end{Highlighting}
\end{Shaded}

\textbf{Scrollback Pager History:}

\begin{Shaded}
\begin{Highlighting}[]
\CommentTok{\# Scrollback separado para pager (MB)}
\DataTypeTok{scrollback\_pager\_history\_size} \DataTypeTok{10}
\end{Highlighting}
\end{Shaded}

\textbf{Fill Window:}

\begin{Shaded}
\begin{Highlighting}[]
\CommentTok{\# Llenar con scrollback al agrandar window}
\DataTypeTok{scrollback\_fill\_enlarged\_window} \DataTypeTok{no}
\end{Highlighting}
\end{Shaded}

\textbf{Scroll Multiplier:}

\begin{Shaded}
\begin{Highlighting}[]
\CommentTok{\# Mouse wheel}
\DataTypeTok{wheel\_scroll\_multiplier} \DataTypeTok{5.0}
\DataTypeTok{wheel\_scroll\_min\_lines} \DataTypeTok{1}

\CommentTok{\# Touch}
\DataTypeTok{touch\_scroll\_multiplier} \DataTypeTok{1.0}
\end{Highlighting}
\end{Shaded}

\subsection{Scrollbar}\label{scrollbar}

\begin{Shaded}
\begin{Highlighting}[]
\CommentTok{\# Cuándo mostrar scrollbar}
\DataTypeTok{scrollbar} \DataTypeTok{scrolled}
\CommentTok{\# Opciones: scrolled, hovered, scrolled{-}and{-}hovered, always, never}

\CommentTok{\# Interactivo}
\DataTypeTok{scrollbar\_interactive} \DataTypeTok{yes}

\CommentTok{\# Jump on click}
\DataTypeTok{scrollbar\_jump\_on\_click} \DataTypeTok{yes}

\CommentTok{\# Ancho}
\DataTypeTok{scrollbar\_width} \DataTypeTok{0.5}
\DataTypeTok{scrollbar\_hover\_width} \DataTypeTok{1}

\CommentTok{\# Opacidad}
\DataTypeTok{scrollbar\_handle\_opacity} \DataTypeTok{0.5}
\DataTypeTok{scrollbar\_track\_opacity} \DataTypeTok{0}

\CommentTok{\# Colores}
\DataTypeTok{scrollbar\_handle\_color} \DataTypeTok{foreground}
\DataTypeTok{scrollbar\_track\_color} \DataTypeTok{foreground}
\end{Highlighting}
\end{Shaded}

\section{Mouse}\label{mouse}

\subsection{Comportamiento Básico}\label{comportamiento-buxe1sico}

\begin{Shaded}
\begin{Highlighting}[]
\CommentTok{\# Ocultar cursor después de inactividad}
\DataTypeTok{mouse\_hide\_wait} \DataTypeTok{3.0}
\CommentTok{\# 0 = nunca}
\CommentTok{\# \textless{}0 = ocultar inmediatamente al escribir}

\CommentTok{\# Focus follows mouse}
\DataTypeTok{focus\_follows\_mouse} \DataTypeTok{no}
\end{Highlighting}
\end{Shaded}

\subsection{URLs}\label{urls}

\begin{Shaded}
\begin{Highlighting}[]
\CommentTok{\# Color y estilo de URLs}
\DataTypeTok{url\_color} \CommentTok{\#0087bd}
\DataTypeTok{url\_style} \DataTypeTok{curly}
\CommentTok{\# Opciones: none, straight, double, curly, dotted, dashed}

\CommentTok{\# Programa para abrir URLs}
\DataTypeTok{open\_url\_with} \DataTypeTok{default}

\CommentTok{\# Prefijos detectados}
\DataTypeTok{url\_prefixes} \DataTypeTok{file} \DataTypeTok{ftp} \DataTypeTok{ftps} \DataTypeTok{gemini} \DataTypeTok{git} \DataTypeTok{gopher} \DataTypeTok{http} \DataTypeTok{https} \DataTypeTok{irc} \DataTypeTok{ircs} \DataTypeTok{kitty} \DataTypeTok{mailto} \DataTypeTok{news} \DataTypeTok{sftp} \DataTypeTok{ssh}

\CommentTok{\# Detectar URLs}
\DataTypeTok{detect\_urls} \DataTypeTok{yes}

\CommentTok{\# Mostrar target de hyperlinks}
\DataTypeTok{show\_hyperlink\_targets} \DataTypeTok{no}

\CommentTok{\# Underline de hyperlinks}
\DataTypeTok{underline\_hyperlinks} \DataTypeTok{hover}
\CommentTok{\# Opciones: hover, always, never}
\end{Highlighting}
\end{Shaded}

\subsection{Selección}\label{selecciuxf3n}

\begin{Shaded}
\begin{Highlighting}[]
\CommentTok{\# Copiar al seleccionar}
\DataTypeTok{copy\_on\_select} \DataTypeTok{no}
\CommentTok{\# Opciones: no, clipboard, a1 (buffer personalizado)}

\CommentTok{\# Limpiar selección al perder clipboard}
\DataTypeTok{clear\_selection\_on\_clipboard\_loss} \DataTypeTok{no}

\CommentTok{\# Caracteres de palabra}
\DataTypeTok{select\_by\_word\_characters} \ErrorTok{@}\DataTypeTok{{-}.}\ErrorTok{/}\DataTypeTok{\_}\ErrorTok{\textasciitilde{}?\&}\OperatorTok{=}\ErrorTok{\%+}\CommentTok{\#}

\CommentTok{\# Intervalo de clicks}
\DataTypeTok{click\_interval} \DataTypeTok{{-}1.0}
\CommentTok{\# {-}1 = usar default del sistema}
\end{Highlighting}
\end{Shaded}

\subsection{Pointer Shape}\label{pointer-shape}

\begin{Shaded}
\begin{Highlighting}[]
\CommentTok{\# Forma del cursor cuando programa agarra mouse}
\DataTypeTok{pointer\_shape\_when\_grabbed} \DataTypeTok{arrow}

\CommentTok{\# Forma default}
\DataTypeTok{default\_pointer\_shape} \DataTypeTok{beam}

\CommentTok{\# Forma al arrastrar}
\DataTypeTok{pointer\_shape\_when\_dragging} \DataTypeTok{beam} \DataTypeTok{crosshair}
\end{Highlighting}
\end{Shaded}

\subsection{Mouse Actions}\label{mouse-actions}

\textbf{Sintaxis:}

\begin{Shaded}
\begin{Highlighting}[]
\DataTypeTok{mouse\_map} \DataTypeTok{button{-}name} \DataTypeTok{event{-}type} \DataTypeTok{modes} \DataTypeTok{action}
\end{Highlighting}
\end{Shaded}

\textbf{Ejemplos:}

\begin{Shaded}
\begin{Highlighting}[]
\CommentTok{\# Click en URL}
\DataTypeTok{mouse\_map} \DataTypeTok{left} \DataTypeTok{click} \DataTypeTok{ungrabbed} \DataTypeTok{mouse\_handle\_click} \DataTypeTok{selection} \DataTypeTok{link} \DataTypeTok{prompt}

\CommentTok{\# Pegar con middle click}
\DataTypeTok{mouse\_map} \DataTypeTok{middle} \DataTypeTok{release} \DataTypeTok{ungrabbed} \DataTypeTok{paste\_from\_selection}

\CommentTok{\# Seleccionar con left press}
\DataTypeTok{mouse\_map} \DataTypeTok{left} \DataTypeTok{press} \DataTypeTok{ungrabbed} \DataTypeTok{mouse\_selection} \DataTypeTok{normal}

\CommentTok{\# Selección rectangular con Ctrl+Alt}
\DataTypeTok{mouse\_map} \DataTypeTok{ctrl}\ErrorTok{+}\DataTypeTok{alt}\ErrorTok{+}\DataTypeTok{left} \DataTypeTok{press} \DataTypeTok{ungrabbed} \DataTypeTok{mouse\_selection} \DataTypeTok{rectangle}

\CommentTok{\# Seleccionar palabra con double click}
\DataTypeTok{mouse\_map} \DataTypeTok{left} \DataTypeTok{doublepress} \DataTypeTok{ungrabbed} \DataTypeTok{mouse\_selection} \DataTypeTok{word}

\CommentTok{\# Seleccionar línea con triple click}
\DataTypeTok{mouse\_map} \DataTypeTok{left} \DataTypeTok{triplepress} \DataTypeTok{ungrabbed} \DataTypeTok{mouse\_selection} \DataTypeTok{line}

\CommentTok{\# Extender selección con right press}
\DataTypeTok{mouse\_map} \DataTypeTok{right} \DataTypeTok{press} \DataTypeTok{ungrabbed} \DataTypeTok{mouse\_selection} \DataTypeTok{extend}
\end{Highlighting}
\end{Shaded}

\textbf{Clearar todos los mouse actions:}

\begin{Shaded}
\begin{Highlighting}[]
\DataTypeTok{clear\_all\_mouse\_actions} \DataTypeTok{yes}
\end{Highlighting}
\end{Shaded}

\section{Kittens (Extensiones)}\label{kittens-extensiones}

\subsection{¿Qué son los Kittens?}\label{quuxe9-son-los-kittens}

Programas auxiliares que extienden funcionalidad de kitty.

\subsection{Kittens Built-in}\label{kittens-built-in}

\subsection{1. icat - Imágenes en
Terminal}\label{icat---imuxe1genes-en-terminal}

\begin{Shaded}
\begin{Highlighting}[]
\CommentTok{\# Mostrar imagen}
\ExtensionTok{kitten}\NormalTok{ icat imagen.png}

\CommentTok{\# Con opciones}
\ExtensionTok{kitten}\NormalTok{ icat }\AttributeTok{{-}{-}align}\NormalTok{ left imagen.png}
\ExtensionTok{kitten}\NormalTok{ icat }\AttributeTok{{-}{-}place}\NormalTok{ 40x40@10x10 imagen.png}
\end{Highlighting}
\end{Shaded}

\textbf{En kitty.conf:}

\begin{Shaded}
\begin{Highlighting}[]
\DataTypeTok{map} \DataTypeTok{f1} \DataTypeTok{launch} \DataTypeTok{{-}{-}stdin{-}source}\OperatorTok{=}\ErrorTok{@}\DataTypeTok{screen\_scrollback} \DataTypeTok{kitten} \DataTypeTok{icat}
\end{Highlighting}
\end{Shaded}

\subsection{2. diff - Diff con
Imágenes}\label{diff---diff-con-imuxe1genes}

\begin{Shaded}
\begin{Highlighting}[]
\CommentTok{\# Diff de archivos}
\ExtensionTok{kitten}\NormalTok{ diff archivo1.txt archivo2.txt}

\CommentTok{\# Con imágenes}
\ExtensionTok{kitten}\NormalTok{ diff imagen1.png imagen2.png}

\CommentTok{\# Diff de directorios}
\ExtensionTok{kitten}\NormalTok{ diff dir1/ dir2/}
\end{Highlighting}
\end{Shaded}

\textbf{Configurar como git difftool:}

\begin{Shaded}
\begin{Highlighting}[]
\FunctionTok{git}\NormalTok{ config }\AttributeTok{{-}{-}global}\NormalTok{ diff.tool kitty}
\FunctionTok{git}\NormalTok{ config }\AttributeTok{{-}{-}global}\NormalTok{ difftool.kitty.cmd }\StringTok{\textquotesingle{}kitten diff $LOCAL $REMOTE\textquotesingle{}}
\end{Highlighting}
\end{Shaded}

\subsection{3. unicode\_input - Input
Unicode}\label{unicode_input---input-unicode}

\begin{Shaded}
\begin{Highlighting}[]
\CommentTok{\# Abrir selector Unicode}
\ExtensionTok{Ctrl+Shift+U}

\CommentTok{\# O}
\ExtensionTok{kitten}\NormalTok{ unicode\_input}
\end{Highlighting}
\end{Shaded}

\subsection{4. themes - Cambiar Temas}\label{themes---cambiar-temas}

\begin{Shaded}
\begin{Highlighting}[]
\CommentTok{\# Browser de temas}
\ExtensionTok{kitten}\NormalTok{ themes}

\CommentTok{\# Listar temas}
\ExtensionTok{kitten}\NormalTok{ themes }\AttributeTok{{-}{-}list}

\CommentTok{\# Cambiar tema}
\ExtensionTok{kitten}\NormalTok{ themes }\AttributeTok{{-}{-}reload{-}in}\OperatorTok{=}\NormalTok{all Dracula}
\end{Highlighting}
\end{Shaded}

\subsection{5. choose-fonts - Elegir
Fuentes}\label{choose-fonts---elegir-fuentes}

\begin{Shaded}
\begin{Highlighting}[]
\CommentTok{\# UI para elegir fuentes}
\ExtensionTok{kitten}\NormalTok{ choose{-}fonts}

\CommentTok{\# Con preview}
\ExtensionTok{kitten}\NormalTok{ choose{-}fonts }\AttributeTok{{-}{-}preview}
\end{Highlighting}
\end{Shaded}

\subsection{6. hints - Selección
Rápida}\label{hints---selecciuxf3n-ruxe1pida}

\begin{Shaded}
\begin{Highlighting}[]
\CommentTok{\# Abrir URLs}
\ExtensionTok{Ctrl+Shift+E}

\CommentTok{\# O personalizado:}
\CommentTok{\# URLs}
\ExtensionTok{kitten}\NormalTok{ hints}

\CommentTok{\# Paths}
\ExtensionTok{kitten}\NormalTok{ hints }\AttributeTok{{-}{-}type}\NormalTok{ path}

\CommentTok{\# Líneas}
\ExtensionTok{kitten}\NormalTok{ hints }\AttributeTok{{-}{-}type}\NormalTok{ line}

\CommentTok{\# Palabras}
\ExtensionTok{kitten}\NormalTok{ hints }\AttributeTok{{-}{-}type}\NormalTok{ word}

\CommentTok{\# Hashes (git)}
\ExtensionTok{kitten}\NormalTok{ hints }\AttributeTok{{-}{-}type}\NormalTok{ hash}

\CommentTok{\# Números de línea}
\ExtensionTok{kitten}\NormalTok{ hints }\AttributeTok{{-}{-}type}\NormalTok{ linenum}

\CommentTok{\# Hyperlinks}
\ExtensionTok{kitten}\NormalTok{ hints }\AttributeTok{{-}{-}type}\NormalTok{ hyperlink}
\end{Highlighting}
\end{Shaded}

\textbf{Configuración:}

\begin{Shaded}
\begin{Highlighting}[]
\CommentTok{\# Abrir URL}
\DataTypeTok{map} \DataTypeTok{ctrl}\ErrorTok{+}\DataTypeTok{shift}\ErrorTok{+}\DataTypeTok{e} \DataTypeTok{open\_url\_with\_hints}

\CommentTok{\# Insertar path}
\DataTypeTok{map} \DataTypeTok{ctrl}\ErrorTok{+}\DataTypeTok{shift}\ErrorTok{+}\DataTypeTok{p}\ErrorTok{\textgreater{}}\DataTypeTok{f} \DataTypeTok{kitten} \DataTypeTok{hints} \DataTypeTok{{-}{-}type} \DataTypeTok{path} \DataTypeTok{{-}{-}program} \DataTypeTok{{-}}

\CommentTok{\# Abrir path}
\DataTypeTok{map} \DataTypeTok{ctrl}\ErrorTok{+}\DataTypeTok{shift}\ErrorTok{+}\DataTypeTok{p}\ErrorTok{\textgreater{}}\DataTypeTok{shift}\ErrorTok{+}\DataTypeTok{f} \DataTypeTok{kitten} \DataTypeTok{hints} \DataTypeTok{{-}{-}type} \DataTypeTok{path}

\CommentTok{\# Insertar archivo (con chooser)}
\DataTypeTok{map} \DataTypeTok{ctrl}\ErrorTok{+}\DataTypeTok{shift}\ErrorTok{+}\DataTypeTok{p}\ErrorTok{\textgreater{}}\DataTypeTok{c} \DataTypeTok{kitten} \DataTypeTok{choose{-}files}

\CommentTok{\# Insertar directorio}
\DataTypeTok{map} \DataTypeTok{ctrl}\ErrorTok{+}\DataTypeTok{shift}\ErrorTok{+}\DataTypeTok{p}\ErrorTok{\textgreater{}}\DataTypeTok{d} \DataTypeTok{kitten} \DataTypeTok{choose{-}files} \DataTypeTok{{-}{-}mode}\OperatorTok{=}\DataTypeTok{dir}

\CommentTok{\# Insertar línea}
\DataTypeTok{map} \DataTypeTok{ctrl}\ErrorTok{+}\DataTypeTok{shift}\ErrorTok{+}\DataTypeTok{p}\ErrorTok{\textgreater{}}\DataTypeTok{l} \DataTypeTok{kitten} \DataTypeTok{hints} \DataTypeTok{{-}{-}type} \DataTypeTok{line} \DataTypeTok{{-}{-}program} \DataTypeTok{{-}}

\CommentTok{\# Insertar palabra}
\DataTypeTok{map} \DataTypeTok{ctrl}\ErrorTok{+}\DataTypeTok{shift}\ErrorTok{+}\DataTypeTok{p}\ErrorTok{\textgreater{}}\DataTypeTok{w} \DataTypeTok{kitten} \DataTypeTok{hints} \DataTypeTok{{-}{-}type} \DataTypeTok{word} \DataTypeTok{{-}{-}program} \DataTypeTok{{-}}

\CommentTok{\# Insertar hash}
\DataTypeTok{map} \DataTypeTok{ctrl}\ErrorTok{+}\DataTypeTok{shift}\ErrorTok{+}\DataTypeTok{p}\ErrorTok{\textgreater{}}\DataTypeTok{h} \DataTypeTok{kitten} \DataTypeTok{hints} \DataTypeTok{{-}{-}type} \DataTypeTok{hash} \DataTypeTok{{-}{-}program} \DataTypeTok{{-}}

\CommentTok{\# Abrir en editor}
\DataTypeTok{map} \DataTypeTok{ctrl}\ErrorTok{+}\DataTypeTok{shift}\ErrorTok{+}\DataTypeTok{p}\ErrorTok{\textgreater{}}\DataTypeTok{n} \DataTypeTok{kitten} \DataTypeTok{hints} \DataTypeTok{{-}{-}type} \DataTypeTok{linenum}

\CommentTok{\# Abrir hyperlink}
\DataTypeTok{map} \DataTypeTok{ctrl}\ErrorTok{+}\DataTypeTok{shift}\ErrorTok{+}\DataTypeTok{p}\ErrorTok{\textgreater{}}\DataTypeTok{y} \DataTypeTok{kitten} \DataTypeTok{hints} \DataTypeTok{{-}{-}type} \DataTypeTok{hyperlink}
\end{Highlighting}
\end{Shaded}

\subsection{7. ssh - SSH Mejorado}\label{ssh---ssh-mejorado}

\begin{Shaded}
\begin{Highlighting}[]
\CommentTok{\# SSH con shell integration}
\ExtensionTok{kitten}\NormalTok{ ssh usuario@servidor}

\CommentTok{\# Con forwarding}
\ExtensionTok{kitten}\NormalTok{ ssh }\AttributeTok{{-}R}\NormalTok{ 52698:localhost:52698 servidor}
\end{Highlighting}
\end{Shaded}

\textbf{Ventajas:} - Shell integration automática - Re-uso de conexiones
- Clonación de config local - Transferencia de terminfo

\subsection{8. transfer - Transferencia de
Archivos}\label{transfer---transferencia-de-archivos}

\begin{Shaded}
\begin{Highlighting}[]
\CommentTok{\# Upload archivo}
\ExtensionTok{kitten}\NormalTok{ transfer archivo.txt servidor:/ruta/}

\CommentTok{\# Download archivo}
\ExtensionTok{kitten}\NormalTok{ transfer servidor:/ruta/archivo.txt ./}

\CommentTok{\# Directorio completo}
\ExtensionTok{kitten}\NormalTok{ transfer }\AttributeTok{{-}{-}recursive}\NormalTok{ directorio/ servidor:/ruta/}
\end{Highlighting}
\end{Shaded}

\subsection{9. clipboard - Clipboard en
SSH}\label{clipboard---clipboard-en-ssh}

\begin{Shaded}
\begin{Highlighting}[]
\CommentTok{\# Copiar a clipboard}
\BuiltInTok{echo} \StringTok{"texto"} \KeywordTok{|} \ExtensionTok{kitten}\NormalTok{ clipboard}

\CommentTok{\# Pegar desde clipboard}
\ExtensionTok{kitten}\NormalTok{ clipboard }\AttributeTok{{-}{-}get{-}clipboard}
\end{Highlighting}
\end{Shaded}

\subsection{10. broadcast - Broadcast
Input}\label{broadcast---broadcast-input}

\begin{Shaded}
\begin{Highlighting}[]
\CommentTok{\# Escribir en múltiples windows}
\ExtensionTok{kitten}\NormalTok{ broadcast}

\CommentTok{\# A windows específicas}
\ExtensionTok{kitten}\NormalTok{ broadcast }\AttributeTok{{-}{-}match{-}tab} \StringTok{"title:server*"}
\end{Highlighting}
\end{Shaded}

\subsection{11. notify - Notificaciones}\label{notify---notificaciones}

\begin{Shaded}
\begin{Highlighting}[]
\CommentTok{\# Enviar notificación}
\ExtensionTok{kitten}\NormalTok{ notify }\StringTok{"Título"} \StringTok{"Mensaje"}

\CommentTok{\# Con icono}
\ExtensionTok{kitten}\NormalTok{ notify }\AttributeTok{{-}{-}icon}\NormalTok{ /path/icon.png }\StringTok{"Título"} \StringTok{"Mensaje"}
\end{Highlighting}
\end{Shaded}

\subsection{12. panel - Panel Desktop}\label{panel---panel-desktop}

\begin{Shaded}
\begin{Highlighting}[]
\CommentTok{\# Crear panel}
\ExtensionTok{kitten}\NormalTok{ panel sh }\AttributeTok{{-}c} \StringTok{\textquotesingle{}while true; do date; sleep 1; done\textquotesingle{}}
\end{Highlighting}
\end{Shaded}

\subsection{13. @ (Remote Control)}\label{remote-control}

\begin{Shaded}
\begin{Highlighting}[]
\CommentTok{\# Ver todos los comandos}
\ExtensionTok{kitten}\NormalTok{ @ }\AttributeTok{{-}{-}help}

\CommentTok{\# Listar windows}
\ExtensionTok{kitten}\NormalTok{ @ ls}

\CommentTok{\# Cambiar colores}
\ExtensionTok{kitten}\NormalTok{ @ set{-}colors foreground=\#ffffff}

\CommentTok{\# Ejecutar comando en window}
\ExtensionTok{kitten}\NormalTok{ @ send{-}text }\AttributeTok{{-}{-}match} \StringTok{"title:vim"} \StringTok{":wq\textbackslash{}n"}
\end{Highlighting}
\end{Shaded}

\subsection{Crear Kittens
Personalizados}\label{crear-kittens-personalizados}

\textbf{Estructura básica:}

\begin{Shaded}
\begin{Highlighting}[]
\CommentTok{\#!/usr/bin/env python3}
\CommentTok{\# mi\_kitten.py}

\ImportTok{from}\NormalTok{ kittens.tui.handler }\ImportTok{import}\NormalTok{ Handler}
\ImportTok{from}\NormalTok{ kitty.key\_encoding }\ImportTok{import}\NormalTok{ KeyEvent}

\KeywordTok{class}\NormalTok{ MyHandler(Handler):}
    \KeywordTok{def}\NormalTok{ initialize(}\VariableTok{self}\NormalTok{):}
        \VariableTok{self}\NormalTok{.}\BuiltInTok{print}\NormalTok{(}\StringTok{"¡Hola desde mi kitten!"}\NormalTok{)}
    
    \KeywordTok{def}\NormalTok{ on\_key(}\VariableTok{self}\NormalTok{, key\_event: KeyEvent):}
        \ControlFlowTok{if}\NormalTok{ key\_event.key }\OperatorTok{==} \StringTok{\textquotesingle{}q\textquotesingle{}}\NormalTok{:}
            \VariableTok{self}\NormalTok{.quit\_loop(}\DecValTok{0}\NormalTok{)}
        \ControlFlowTok{return} \VariableTok{True}

\KeywordTok{def}\NormalTok{ main(args):}
\NormalTok{    handler }\OperatorTok{=}\NormalTok{ MyHandler()}
\NormalTok{    handler.run()}

\ControlFlowTok{if} \VariableTok{\_\_name\_\_} \OperatorTok{==} \StringTok{\textquotesingle{}\_\_main\_\_\textquotesingle{}}\NormalTok{:}
    \ImportTok{import}\NormalTok{ sys}
\NormalTok{    main(sys.argv[}\DecValTok{1}\NormalTok{:])}
\end{Highlighting}
\end{Shaded}

\textbf{Usar:}

\begin{Shaded}
\begin{Highlighting}[]
\ExtensionTok{kitten}\NormalTok{ @ run{-}kitten /path/to/mi\_kitten.py}
\end{Highlighting}
\end{Shaded}

\section{Control Remoto}\label{control-remoto}

\subsection{Habilitar Remote Control}\label{habilitar-remote-control}

\begin{Shaded}
\begin{Highlighting}[]
\CommentTok{\# En kitty.conf}
\DataTypeTok{allow\_remote\_control} \DataTypeTok{yes}

\CommentTok{\# Socket Unix}
\DataTypeTok{listen\_on} \DataTypeTok{unix}\ErrorTok{:/}\DataTypeTok{tmp}\ErrorTok{/}\DataTypeTok{kitty}

\CommentTok{\# O TCP}
\DataTypeTok{listen\_on} \DataTypeTok{tcp}\ErrorTok{:}\DataTypeTok{localhost}\ErrorTok{:}\DataTypeTok{12345}
\end{Highlighting}
\end{Shaded}

\subsection{Uso Básico}\label{uso-buxe1sico}

\begin{Shaded}
\begin{Highlighting}[]
\CommentTok{\# Listar windows}
\ExtensionTok{kitty}\NormalTok{ @ ls}

\CommentTok{\# Crear nueva window}
\ExtensionTok{kitty}\NormalTok{ @ launch}

\CommentTok{\# Cerrar window}
\ExtensionTok{kitty}\NormalTok{ @ close{-}window}

\CommentTok{\# Cambiar título}
\ExtensionTok{kitty}\NormalTok{ @ set{-}window{-}title }\StringTok{"Nuevo Título"}

\CommentTok{\# Enviar texto}
\ExtensionTok{kitty}\NormalTok{ @ send{-}text }\StringTok{"echo hello\textbackslash{}n"}
\end{Highlighting}
\end{Shaded}

\subsection{Comandos Principales}\label{comandos-principales}

\textbf{Información:}

\begin{Shaded}
\begin{Highlighting}[]
\ExtensionTok{kitty}\NormalTok{ @ ls                    }\CommentTok{\# Listar todo}
\ExtensionTok{kitty}\NormalTok{ @ get{-}colors            }\CommentTok{\# Ver colores actuales}
\ExtensionTok{kitty}\NormalTok{ @ get{-}text              }\CommentTok{\# Obtener texto de window}
\end{Highlighting}
\end{Shaded}

\textbf{Crear/Cerrar:}

\begin{Shaded}
\begin{Highlighting}[]
\ExtensionTok{kitty}\NormalTok{ @ launch                }\CommentTok{\# Nueva window}
\ExtensionTok{kitty}\NormalTok{ @ launch }\AttributeTok{{-}{-}type}\OperatorTok{=}\NormalTok{tab     }\CommentTok{\# Nueva tab}
\ExtensionTok{kitty}\NormalTok{ @ launch }\AttributeTok{{-}{-}type}\OperatorTok{=}\NormalTok{os{-}window  }\CommentTok{\# Nueva OS window}
\ExtensionTok{kitty}\NormalTok{ @ close{-}window}
\ExtensionTok{kitty}\NormalTok{ @ close{-}tab}
\end{Highlighting}
\end{Shaded}

\textbf{Modificar:}

\begin{Shaded}
\begin{Highlighting}[]
\ExtensionTok{kitty}\NormalTok{ @ set{-}colors foreground=\#ff0000}
\ExtensionTok{kitty}\NormalTok{ @ set{-}window{-}title }\StringTok{"Título"}
\ExtensionTok{kitty}\NormalTok{ @ set{-}tab{-}title }\StringTok{"Tab"}
\ExtensionTok{kitty}\NormalTok{ @ resize{-}window }\AttributeTok{{-}{-}increment}\NormalTok{ 10}
\ExtensionTok{kitty}\NormalTok{ @ goto{-}layout tall}
\end{Highlighting}
\end{Shaded}

\textbf{Enviar Input:}

\begin{Shaded}
\begin{Highlighting}[]
\ExtensionTok{kitty}\NormalTok{ @ send{-}text }\StringTok{"texto"}
\ExtensionTok{kitty}\NormalTok{ @ send{-}text }\AttributeTok{{-}{-}match} \StringTok{"title:vim"} \StringTok{":wq\textbackslash{}n"}
\end{Highlighting}
\end{Shaded}

\textbf{Focus:}

\begin{Shaded}
\begin{Highlighting}[]
\ExtensionTok{kitty}\NormalTok{ @ focus{-}window }\AttributeTok{{-}{-}match} \StringTok{"title:editor"}
\ExtensionTok{kitty}\NormalTok{ @ focus{-}tab }\AttributeTok{{-}{-}match} \StringTok{"index:2"}
\end{Highlighting}
\end{Shaded}

\subsection{Matching}\label{matching}

Usar \texttt{-\/-match} para seleccionar windows/tabs:

\begin{Shaded}
\begin{Highlighting}[]
\CommentTok{\# Por título}
\ExtensionTok{{-}{-}match} \StringTok{"title:editor"}

\CommentTok{\# Por ID}
\ExtensionTok{{-}{-}match} \StringTok{"id:42"}

\CommentTok{\# Por estado}
\ExtensionTok{{-}{-}match} \StringTok{"state:focused"}

\CommentTok{\# Múltiples condiciones}
\ExtensionTok{{-}{-}match} \StringTok{"title:vim and state:focused"}
\end{Highlighting}
\end{Shaded}

\subsection{Remote Control desde
Mapeos}\label{remote-control-desde-mapeos}

\begin{Shaded}
\begin{Highlighting}[]
\CommentTok{\# Crear window con comando específico}
\DataTypeTok{map} \DataTypeTok{f1} \DataTypeTok{remote\_control} \DataTypeTok{launch} \DataTypeTok{{-}{-}type}\OperatorTok{=}\DataTypeTok{tab} \DataTypeTok{vim}

\CommentTok{\# Cambiar colores}
\DataTypeTok{map} \DataTypeTok{f2} \DataTypeTok{remote\_control} \DataTypeTok{set{-}colors} \DataTypeTok{{-}{-}configured} \ErrorTok{\textasciitilde{}/}\DataTypeTok{colors.conf}

\CommentTok{\# Broadcast}
\DataTypeTok{map} \DataTypeTok{f3} \DataTypeTok{remote\_control} \DataTypeTok{send{-}text} \DataTypeTok{{-}{-}all} \DataTypeTok{"clear}\SpecialCharTok{\textbackslash{}n}\DataTypeTok{"}
\end{Highlighting}
\end{Shaded}

\subsection{Remote Control via SSH}\label{remote-control-via-ssh}

\begin{Shaded}
\begin{Highlighting}[]
\CommentTok{\# Desde máquina remota}
\BuiltInTok{export} \VariableTok{KITTY\_REMOTE\_CONTROL\_PASSWORD}\OperatorTok{=}\StringTok{"mi\_password"}
\ExtensionTok{kitty}\NormalTok{ @ }\AttributeTok{{-}{-}password}\OperatorTok{=}\VariableTok{$KITTY\_REMOTE\_CONTROL\_PASSWORD}\NormalTok{ ls}
\end{Highlighting}
\end{Shaded}

\textbf{Con kitten ssh:}

\begin{Shaded}
\begin{Highlighting}[]
\CommentTok{\# Automáticamente habilita remote control}
\ExtensionTok{kitten}\NormalTok{ ssh servidor}

\CommentTok{\# Desde el servidor}
\ExtensionTok{kitty}\NormalTok{ @ ls  }\CommentTok{\# ¡Funciona!}
\end{Highlighting}
\end{Shaded}

\section{Sesiones}\label{sesiones}

\subsection{¿Qué son las Sesiones?}\label{quuxe9-son-las-sesiones}

Archivos que definen:

\begin{itemize}
\tightlist
\item
  Layout de tabs y windows
\item
  Directorio de trabajo
\item
  Comandos a ejecutar
\item
  Configuración específica
\end{itemize}

\subsection{Crear Sesión}\label{crear-sesiuxf3n}

\textbf{Archivo de sesión:}

\begin{verbatim}
# mi_sesion.kitty

# Nueva ventana
new_window
cd ~/proyectos/app
launch vim

# Nueva ventana
new_window
cd ~/proyectos/app
launch

# Nueva tab
new_tab editor
cd ~/proyectos/app
layout tall
launch vim src/main.py
launch

# Nueva tab con split
new_tab servers
launch ssh server1
launch --location=hsplit ssh server2
\end{verbatim}

\subsection{Usar Sesión}\label{usar-sesiuxf3n}

\textbf{Al iniciar kitty:}

\begin{Shaded}
\begin{Highlighting}[]
\ExtensionTok{kitty} \AttributeTok{{-}{-}session}\NormalTok{ mi\_sesion.kitty}
\end{Highlighting}
\end{Shaded}

\textbf{Desde configuración:}

\begin{Shaded}
\begin{Highlighting}[]
\DataTypeTok{startup\_session} \DataTypeTok{mi\_sesion.kitty}
\end{Highlighting}
\end{Shaded}

\textbf{Desde atajo:}

\begin{Shaded}
\begin{Highlighting}[]
\DataTypeTok{map} \DataTypeTok{f1} \DataTypeTok{launch} \DataTypeTok{{-}{-}type}\OperatorTok{=}\DataTypeTok{os{-}window} \DataTypeTok{{-}{-}session}\OperatorTok{=}\DataTypeTok{mi\_sesion.kitty}
\end{Highlighting}
\end{Shaded}

\subsection{Comandos de Sesión}\label{comandos-de-sesiuxf3n}

\textbf{new\_tab:}

\begin{verbatim}
new_tab [nombre]
cd /directorio
layout tall
launch comando
\end{verbatim}

\textbf{new\_window:}

\begin{verbatim}
new_window
cd /directorio
launch comando
\end{verbatim}

\textbf{launch:}

\begin{verbatim}
launch [--location=split] comando
\end{verbatim}

\textbf{cd:}

\begin{verbatim}
cd /ruta/absoluta
cd ~/relativa
\end{verbatim}

\textbf{layout:}

\begin{verbatim}
layout tall
layout stack
layout splits
\end{verbatim}

\textbf{focus:}

\begin{verbatim}
focus
\end{verbatim}

\textbf{enabled\_layouts:}

\begin{verbatim}
enabled_layouts tall,stack,grid
\end{verbatim}

\subsection{Ejemplo Completo}\label{ejemplo-completo}

\begin{verbatim}
# Sesión de desarrollo web
# ~/proyectos/web_dev.kitty

# Tab principal - Editor
new_tab Editor
cd ~/proyectos/webapp
layout tall:bias=70;full_size=1
launch nvim .

# Tab de servidores
new_tab Servers
cd ~/proyectos/webapp
layout splits
launch npm run dev
launch --location=hsplit npm run tailwind

# Tab de terminal
new_tab Terminal
cd ~/proyectos/webapp
launch

# Tab de git
new_tab Git
cd ~/proyectos/webapp
launch lazygit

# Focus en primera tab
focus_tab 1
\end{verbatim}

\subsection{Guardar Sesión Actual}\label{guardar-sesiuxf3n-actual}

Kitty no tiene comando built-in para guardar, pero puedes usar remote
control:

\begin{Shaded}
\begin{Highlighting}[]
\CommentTok{\# Script para guardar sesión}
\ExtensionTok{kitty}\NormalTok{ @ ls }\KeywordTok{|} \ExtensionTok{python3}\NormalTok{ crear\_sesion.py }\OperatorTok{\textgreater{}}\NormalTok{ sesion\_actual.kitty}
\end{Highlighting}
\end{Shaded}

\subsection{Switching Sessions}\label{switching-sessions}

\textbf{Con goto\_session:}

\begin{Shaded}
\begin{Highlighting}[]
\CommentTok{\# Crear/cambiar a sesión}
\DataTypeTok{map} \DataTypeTok{f1} \DataTypeTok{goto\_session} \DataTypeTok{dev}
\DataTypeTok{map} \DataTypeTok{f2} \DataTypeTok{goto\_session} \DataTypeTok{personal}
\end{Highlighting}
\end{Shaded}

\textbf{Cerrar sesión:}

\begin{Shaded}
\begin{Highlighting}[]
\CommentTok{\# Cerrar todas las windows de una sesión}
\DataTypeTok{map} \DataTypeTok{f3} \DataTypeTok{close\_session} \DataTypeTok{dev}
\end{Highlighting}
\end{Shaded}

\section{Integración con Shell}\label{integraciuxf3n-con-shell}

\subsection{Shell Integration}\label{shell-integration}

Kitty puede integrarse con bash, zsh y fish para:

\begin{itemize}
\tightlist
\item
  Saltar a prompts anteriores
\item
  Ver output de comandos
\item
  Tracking de directorios
\item
  Indicadores visuales
\end{itemize}

\subsection{Habilitar}\label{habilitar}

\begin{Shaded}
\begin{Highlighting}[]
\DataTypeTok{shell\_integration} \DataTypeTok{enabled}
\end{Highlighting}
\end{Shaded}

\textbf{Deshabilitar características específicas:}

\begin{Shaded}
\begin{Highlighting}[]
\CommentTok{\# Deshabilitar cursor shape en prompt}
\DataTypeTok{shell\_integration} \DataTypeTok{no{-}cursor}

\CommentTok{\# Deshabilitar título automático}
\DataTypeTok{shell\_integration} \DataTypeTok{no{-}title}

\CommentTok{\# Deshabilitar tracking de cwd}
\DataTypeTok{shell\_integration} \DataTypeTok{no{-}cwd}

\CommentTok{\# Deshabilitar prompt marks}
\DataTypeTok{shell\_integration} \DataTypeTok{no{-}prompt{-}mark}

\CommentTok{\# Deshabilitar completion}
\DataTypeTok{shell\_integration} \DataTypeTok{no{-}complete}

\CommentTok{\# Deshabilitar sudo}
\DataTypeTok{shell\_integration} \DataTypeTok{no{-}sudo}
\end{Highlighting}
\end{Shaded}

\subsection{Configuración Manual}\label{configuraciuxf3n-manual}

Si la integración automática no funciona:

\textbf{Bash (\textasciitilde/.bashrc):}

\begin{Shaded}
\begin{Highlighting}[]
\ControlFlowTok{if} \KeywordTok{[[} \OtherTok{{-}n} \StringTok{"}\VariableTok{$KITTY\_INSTALLATION\_DIR}\StringTok{"} \KeywordTok{]];} \ControlFlowTok{then}
    \BuiltInTok{export} \VariableTok{KITTY\_SHELL\_INTEGRATION}\OperatorTok{=}\StringTok{"enabled"}
    \BuiltInTok{source} \StringTok{"}\VariableTok{$KITTY\_INSTALLATION\_DIR}\StringTok{/shell{-}integration/bash/kitty.bash"}
\ControlFlowTok{fi}
\end{Highlighting}
\end{Shaded}

\textbf{Zsh (\textasciitilde/.zshrc):}

\begin{Shaded}
\begin{Highlighting}[]
\ControlFlowTok{if} \KeywordTok{[[} \OtherTok{{-}n} \StringTok{"}\VariableTok{$KITTY\_INSTALLATION\_DIR}\StringTok{"} \KeywordTok{]];} \ControlFlowTok{then}
    \BuiltInTok{export} \VariableTok{KITTY\_SHELL\_INTEGRATION}\OperatorTok{=}\StringTok{"enabled"}
    \BuiltInTok{source} \StringTok{"}\VariableTok{$KITTY\_INSTALLATION\_DIR}\StringTok{/shell{-}integration/zsh/kitty.zsh"}
\ControlFlowTok{fi}
\end{Highlighting}
\end{Shaded}

\textbf{Fish (\textasciitilde/.config/fish/config.fish):}

\begin{Shaded}
\begin{Highlighting}[]
\NormalTok{if set {-}q KITTY\_INSTALLATION\_DIR}
\NormalTok{    set {-}{-}global KITTY\_SHELL\_INTEGRATION enabled}
\NormalTok{    source "$KITTY\_INSTALLATION\_DIR/shell{-}integration/fish/vendor\_conf.d/kitty{-}shell{-}integration.fish"}
\NormalTok{    set {-}{-}prepend fish\_complete\_path "$KITTY\_INSTALLATION\_DIR/shell{-}integration/fish/vendor\_completions.d"}
\NormalTok{end}
\end{Highlighting}
\end{Shaded}

\subsection{Características con Shell
Integration}\label{caracteruxedsticas-con-shell-integration}

\textbf{Saltar a prompts:}

\begin{Shaded}
\begin{Highlighting}[]
\ExtensionTok{Ctrl+Shift+Z}  \CommentTok{\# Prompt anterior}
\ExtensionTok{Ctrl+Shift+X}  \CommentTok{\# Prompt siguiente}
\end{Highlighting}
\end{Shaded}

\textbf{Ver output de comandos:}

\begin{Shaded}
\begin{Highlighting}[]
\ExtensionTok{Ctrl+Shift+G}  \CommentTok{\# Ver último comando en pager}
\end{Highlighting}
\end{Shaded}

\textbf{Mouse:}

\begin{itemize}
\tightlist
\item
  Click en prompt para mover cursor
\item
  Right-click en comando para ver output
\item
  Click en error para copiar
\end{itemize}

\textbf{Indicadores visuales:}

\begin{itemize}
\tightlist
\item
  Marca en prompt
\item
  Color diferente para comandos exitosos/fallidos
\end{itemize}

\section{Marcas (Marks)}\label{marcas-marks}

\subsection{¿Qué son las Marcas?}\label{quuxe9-son-las-marcas}

Sistema para resaltar texto basado en regex.

\subsection{Tipos de Marcas}\label{tipos-de-marcas}

Hay 3 tipos (cada uno con color diferente):

\begin{itemize}
\tightlist
\item
  mark1 (light steel blue)
\item
  mark2 (beige)
\item
  mark3 (violet)
\end{itemize}

\subsection{Configurar Colores}\label{configurar-colores}

\begin{Shaded}
\begin{Highlighting}[]
\CommentTok{\# Mark tipo 1}
\DataTypeTok{mark1\_foreground} \DataTypeTok{black}
\DataTypeTok{mark1\_background} \CommentTok{\#98d3cb}

\CommentTok{\# Mark tipo 2}
\DataTypeTok{mark2\_foreground} \DataTypeTok{black}
\DataTypeTok{mark2\_background} \CommentTok{\#f2dcd3}

\CommentTok{\# Mark tipo 3}
\DataTypeTok{mark3\_foreground} \DataTypeTok{black}
\DataTypeTok{mark3\_background} \CommentTok{\#f274bc}
\end{Highlighting}
\end{Shaded}

\subsection{Crear Marcas}\label{crear-marcas}

\textbf{Con regex:}

\begin{Shaded}
\begin{Highlighting}[]
\CommentTok{\# Marcar URLs}
\DataTypeTok{map} \DataTypeTok{f1} \DataTypeTok{create\_marker} \DataTypeTok{text} \DataTypeTok{1} \DataTypeTok{https}\ErrorTok{?://}\KeywordTok{[}\ErrorTok{\^{}\textbackslash{}}\DataTypeTok{s}\ErrorTok{]+}

\CommentTok{\# Marcar errores}
\DataTypeTok{map} \DataTypeTok{f2} \DataTypeTok{create\_marker} \DataTypeTok{text} \DataTypeTok{1} \DataTypeTok{ERROR}\ErrorTok{|}\DataTypeTok{FATAL}

\CommentTok{\# Marcar }\AlertTok{TODO}
\DataTypeTok{map} \DataTypeTok{f3} \DataTypeTok{create\_marker} \DataTypeTok{text} \DataTypeTok{2} \DataTypeTok{TODO}\ErrorTok{|}\DataTypeTok{FIXME}\ErrorTok{|}\DataTypeTok{XXX}
\end{Highlighting}
\end{Shaded}

\textbf{Desde línea de comando:}

\begin{Shaded}
\begin{Highlighting}[]
\CommentTok{\# Marcar palabra}
\ExtensionTok{kitty}\NormalTok{ @ create{-}marker text 1 }\StringTok{"error"}

\CommentTok{\# Marcar líneas que contengan}
\ExtensionTok{kitty}\NormalTok{ @ create{-}marker iregex 1 }\StringTok{"\^{}error.*"}
\end{Highlighting}
\end{Shaded}

\subsection{Usar Marcas}\label{usar-marcas}

\textbf{Toggle:}

\begin{Shaded}
\begin{Highlighting}[]
\DataTypeTok{map} \DataTypeTok{f1} \DataTypeTok{toggle\_marker} \DataTypeTok{text} \DataTypeTok{1} \DataTypeTok{error}
\end{Highlighting}
\end{Shaded}

\textbf{Saltar:}

\begin{Shaded}
\begin{Highlighting}[]
\CommentTok{\# Siguiente marca}
\DataTypeTok{map} \DataTypeTok{f2} \DataTypeTok{scroll\_to\_mark} \DataTypeTok{next} \DataTypeTok{1}

\CommentTok{\# Anterior marca}
\DataTypeTok{map} \DataTypeTok{f3} \DataTypeTok{scroll\_to\_mark} \DataTypeTok{prev} \DataTypeTok{1}

\CommentTok{\# Cualquier tipo}
\DataTypeTok{map} \DataTypeTok{f4} \DataTypeTok{scroll\_to\_mark} \DataTypeTok{next} \DataTypeTok{0}
\end{Highlighting}
\end{Shaded}

\textbf{Eliminar:}

\begin{Shaded}
\begin{Highlighting}[]
\DataTypeTok{map} \DataTypeTok{f5} \DataTypeTok{remove\_marker}
\end{Highlighting}
\end{Shaded}

\subsection{Ejemplo Completo}\label{ejemplo-completo-1}

\begin{Shaded}
\begin{Highlighting}[]
\CommentTok{\# Marcar diferentes niveles de log}
\DataTypeTok{map} \DataTypeTok{ctrl}\ErrorTok{+}\DataTypeTok{shift}\ErrorTok{+}\DataTypeTok{1} \DataTypeTok{create\_marker} \DataTypeTok{iregex} \DataTypeTok{1} \DataTypeTok{"\^{}.*ERROR.*$"}
\DataTypeTok{map} \DataTypeTok{ctrl}\ErrorTok{+}\DataTypeTok{shift}\ErrorTok{+}\DataTypeTok{2} \DataTypeTok{create\_marker} \DataTypeTok{iregex} \DataTypeTok{2} \DataTypeTok{"\^{}.*WARNING.*$"}
\DataTypeTok{map} \DataTypeTok{ctrl}\ErrorTok{+}\DataTypeTok{shift}\ErrorTok{+}\DataTypeTok{3} \DataTypeTok{create\_marker} \DataTypeTok{iregex} \DataTypeTok{3} \DataTypeTok{"\^{}.*INFO.*$"}

\CommentTok{\# Navegar entre marcas}
\DataTypeTok{map} \DataTypeTok{ctrl}\ErrorTok{+}\DataTypeTok{shift}\ErrorTok{+}\DataTypeTok{n} \DataTypeTok{scroll\_to\_mark} \DataTypeTok{next}
\DataTypeTok{map} \DataTypeTok{ctrl}\ErrorTok{+}\DataTypeTok{shift}\ErrorTok{+}\DataTypeTok{p} \DataTypeTok{scroll\_to\_mark} \DataTypeTok{prev}

\CommentTok{\# Toggle marcas}
\DataTypeTok{map} \DataTypeTok{ctrl}\ErrorTok{+}\DataTypeTok{shift}\ErrorTok{+}\DataTypeTok{m} \DataTypeTok{toggle\_marker} \DataTypeTok{iregex} \DataTypeTok{1} \DataTypeTok{"\^{}.*ERROR.*$"}

\CommentTok{\# Limpiar todas las marcas}
\DataTypeTok{map} \DataTypeTok{ctrl}\ErrorTok{+}\DataTypeTok{shift}\ErrorTok{+}\DataTypeTok{c} \DataTypeTok{remove\_marker}
\end{Highlighting}
\end{Shaded}

\section{Características Avanzadas}\label{caracteruxedsticas-avanzadas}

\subsection{Multiple Copy/Paste
Buffers}\label{multiple-copypaste-buffers}

\begin{Shaded}
\begin{Highlighting}[]
\CommentTok{\# Definir buffer \textquotesingle{}a\textquotesingle{}}
\DataTypeTok{map} \DataTypeTok{f1} \DataTypeTok{copy\_to\_buffer} \DataTypeTok{a}
\DataTypeTok{map} \DataTypeTok{f2} \DataTypeTok{paste\_from\_buffer} \DataTypeTok{a}

\CommentTok{\# Definir buffer \textquotesingle{}b\textquotesingle{}}
\DataTypeTok{map} \DataTypeTok{f3} \DataTypeTok{copy\_to\_buffer} \DataTypeTok{b}
\DataTypeTok{map} \DataTypeTok{f4} \DataTypeTok{paste\_from\_buffer} \DataTypeTok{b}
\end{Highlighting}
\end{Shaded}

\subsection{Keyboard Modes}\label{keyboard-modes}

Modos de teclado para mapeos contextuales:

\begin{Shaded}
\begin{Highlighting}[]
\CommentTok{\# Crear modo}
\DataTypeTok{map} \DataTypeTok{f1} \DataTypeTok{push\_keyboard\_mode} \DataTypeTok{resize}

\CommentTok{\# En modo resize}
\DataTypeTok{map} \DataTypeTok{{-}{-}mode} \DataTypeTok{resize} \DataTypeTok{h} \DataTypeTok{resize\_window} \DataTypeTok{narrower}
\DataTypeTok{map} \DataTypeTok{{-}{-}mode} \DataTypeTok{resize} \DataTypeTok{l} \DataTypeTok{resize\_window} \DataTypeTok{wider}
\DataTypeTok{map} \DataTypeTok{{-}{-}mode} \DataTypeTok{resize} \DataTypeTok{k} \DataTypeTok{resize\_window} \DataTypeTok{taller}
\DataTypeTok{map} \DataTypeTok{{-}{-}mode} \DataTypeTok{resize} \DataTypeTok{j} \DataTypeTok{resize\_window} \DataTypeTok{shorter}
\DataTypeTok{map} \DataTypeTok{{-}{-}mode} \DataTypeTok{resize} \DataTypeTok{esc} \DataTypeTok{pop\_keyboard\_mode}
\end{Highlighting}
\end{Shaded}

\subsection{Per-Application Mappings}\label{per-application-mappings}

\begin{Shaded}
\begin{Highlighting}[]
\CommentTok{\# Solo para vim}
\DataTypeTok{map} \DataTypeTok{{-}{-}when{-}focus{-}on} \DataTypeTok{vim} \DataTypeTok{ctrl}\ErrorTok{+}\DataTypeTok{j} \DataTypeTok{send\_text} \ErrorTok{\textbackslash{}}\DataTypeTok{x0a}

\CommentTok{\# Solo para less}
\DataTypeTok{map} \DataTypeTok{{-}{-}when{-}focus{-}on} \DataTypeTok{less} \DataTypeTok{j} \DataTypeTok{send\_text} \DataTypeTok{j}
\end{Highlighting}
\end{Shaded}

\subsection{Environment Variables}\label{environment-variables}

\begin{Shaded}
\begin{Highlighting}[]
\CommentTok{\# Establecer variables}
\DataTypeTok{env} \DataTypeTok{VAR}\OperatorTok{=}\DataTypeTok{value}
\DataTypeTok{env} \DataTypeTok{PATH}\OperatorTok{=}\ErrorTok{/}\DataTypeTok{custom}\ErrorTok{/}\DataTypeTok{path}\ErrorTok{:$\{}\DataTypeTok{PATH}\ErrorTok{\}}

\CommentTok{\# Leer de shell}
\DataTypeTok{env} \DataTypeTok{read\_from\_shell}\OperatorTok{=}\DataTypeTok{PATH} \DataTypeTok{LANG} \DataTypeTok{LC\_}\ErrorTok{*} \DataTypeTok{XDG\_}\ErrorTok{*}
\end{Highlighting}
\end{Shaded}

\subsection{Watchers}\label{watchers}

Scripts Python que reaccionan a eventos:

\textbf{\textasciitilde/.config/kitty/watchers.py:}

\begin{Shaded}
\begin{Highlighting}[]
\ImportTok{from}\NormalTok{ kitty.boss }\ImportTok{import}\NormalTok{ Boss}

\KeywordTok{def}\NormalTok{ on\_focus\_change(boss: Boss, window, data):}
    \ControlFlowTok{if}\NormalTok{ data[}\StringTok{\textquotesingle{}focused\textquotesingle{}}\NormalTok{]:}
        \CommentTok{\# Window ganó focus}
        \ControlFlowTok{pass}
    \ControlFlowTok{else}\NormalTok{:}
        \CommentTok{\# Window perdió focus}
        \ControlFlowTok{pass}

\KeywordTok{def}\NormalTok{ on\_close(boss: Boss, window, data):}
    \CommentTok{\# Window está cerrándose}
    \ControlFlowTok{pass}

\KeywordTok{def}\NormalTok{ on\_resize(boss: Boss, window, data):}
    \CommentTok{\# Window fue redimensionada}
    \ControlFlowTok{pass}

\KeywordTok{def}\NormalTok{ on\_title\_change(boss: Boss, window, data):}
\NormalTok{    new\_title }\OperatorTok{=}\NormalTok{ data[}\StringTok{\textquotesingle{}title\textquotesingle{}}\NormalTok{]}
    \CommentTok{\# Hacer algo con nuevo título}
    \ControlFlowTok{pass}
\end{Highlighting}
\end{Shaded}

\textbf{Habilitar:}

\begin{Shaded}
\begin{Highlighting}[]
\DataTypeTok{watcher} \ErrorTok{\textasciitilde{}/}\DataTypeTok{.config}\ErrorTok{/}\DataTypeTok{kitty}\ErrorTok{/}\DataTypeTok{watchers.py}
\end{Highlighting}
\end{Shaded}

\subsection{Filtros de Notificación}\label{filtros-de-notificaciuxf3n}

\begin{Shaded}
\begin{Highlighting}[]
\CommentTok{\# Filtrar notificaciones por campo}
\DataTypeTok{filter\_notification} \DataTypeTok{title}\ErrorTok{:}\DataTypeTok{hello} \DataTypeTok{or} \DataTypeTok{body}\ErrorTok{:}\DataTypeTok{"abc.*def"}
\DataTypeTok{filter\_notification} \DataTypeTok{app}\ErrorTok{:}\DataTypeTok{"[ng]?vim"} \DataTypeTok{and} \DataTypeTok{not} \DataTypeTok{body}\ErrorTok{:}\DataTypeTok{"(?i)update"}

\CommentTok{\# Filtrar todas}
\DataTypeTok{filter\_notification} \DataTypeTok{all}
\end{Highlighting}
\end{Shaded}

\subsection{Allow Cloning}\label{allow-cloning}

\begin{Shaded}
\begin{Highlighting}[]
\CommentTok{\# Permitir clone{-}in{-}kitty}
\DataTypeTok{allow\_cloning} \DataTypeTok{ask}
\CommentTok{\# Opciones: yes, no, ask}
\end{Highlighting}
\end{Shaded}

\subsection{Notificaciones de Comando}\label{notificaciones-de-comando}

\begin{Shaded}
\begin{Highlighting}[]
\CommentTok{\# Notificar cuando comando termina}
\DataTypeTok{notify\_on\_cmd\_finish} \DataTypeTok{unfocused} \DataTypeTok{10.0}

\CommentTok{\# Con bell}
\DataTypeTok{notify\_on\_cmd\_finish} \DataTypeTok{invisible} \DataTypeTok{10.0} \DataTypeTok{bell}

\CommentTok{\# Con comando personalizado}
\DataTypeTok{notify\_on\_cmd\_finish} \DataTypeTok{invisible} \DataTypeTok{10.0} \DataTypeTok{command} \DataTypeTok{notify{-}send} \DataTypeTok{"Done"} \DataTypeTok{"\%c"}
\end{Highlighting}
\end{Shaded}

\section{Optimización de
Rendimiento}\label{optimizaciuxf3n-de-rendimiento}

\subsection{Repaint Delay}\label{repaint-delay}

\begin{Shaded}
\begin{Highlighting}[]
\CommentTok{\# Delay entre actualizaciones (ms)}
\DataTypeTok{repaint\_delay} \DataTypeTok{10}
\CommentTok{\# Menor = más FPS pero más CPU}
\end{Highlighting}
\end{Shaded}

\subsection{Input Delay}\label{input-delay}

\begin{Shaded}
\begin{Highlighting}[]
\CommentTok{\# Delay antes de procesar input (ms)}
\DataTypeTok{input\_delay} \DataTypeTok{3}
\CommentTok{\# Menor = más responsive pero más CPU}
\end{Highlighting}
\end{Shaded}

\subsection{Sync to Monitor}\label{sync-to-monitor}

\begin{Shaded}
\begin{Highlighting}[]
\CommentTok{\# Sincronizar con refresh rate}
\DataTypeTok{sync\_to\_monitor} \DataTypeTok{yes}
\CommentTok{\# no = más FPS pero posible tearing}
\end{Highlighting}
\end{Shaded}

\subsection{Scrollback}\label{scrollback-1}

\begin{Shaded}
\begin{Highlighting}[]
\CommentTok{\# No mantener scrollback enorme}
\DataTypeTok{scrollback\_lines} \DataTypeTok{10000}

\CommentTok{\# Usar scrollback\_pager\_history\_size}
\DataTypeTok{scrollback\_pager\_history\_size} \DataTypeTok{10}
\end{Highlighting}
\end{Shaded}

\subsection{Text Composition}\label{text-composition}

\begin{Shaded}
\begin{Highlighting}[]
\CommentTok{\# Estrategia de composición}
\DataTypeTok{text\_composition\_strategy} \DataTypeTok{platform}
\CommentTok{\# platform, legacy, o números para ajuste fino}
\end{Highlighting}
\end{Shaded}

\subsection{OpenGL}\label{opengl}

\begin{Shaded}
\begin{Highlighting}[]
\CommentTok{\# Verificar que usa GPU}
\ExtensionTok{glxinfo} \KeywordTok{|} \FunctionTok{grep} \StringTok{"OpenGL renderer"}

\CommentTok{\# Actualizar drivers GPU}
\FunctionTok{sudo}\NormalTok{ apt update }\KeywordTok{\&\&} \FunctionTok{sudo}\NormalTok{ apt upgrade}
\end{Highlighting}
\end{Shaded}

\section{Trucos y Tips}\label{trucos-y-tips}

\subsection{Cambiar Directorio del Shell al
Salir}\label{cambiar-directorio-del-shell-al-salir}

\textbf{Bash/Zsh:}

\begin{Shaded}
\begin{Highlighting}[]
\CommentTok{\# Añadir a \textasciitilde{}/.bashrc o \textasciitilde{}/.zshrc}
\BuiltInTok{alias}\NormalTok{ icat=}\StringTok{"kitten icat"}
\end{Highlighting}
\end{Shaded}

\subsection{Previsualizar Imágenes}\label{previsualizar-imuxe1genes}

\begin{Shaded}
\begin{Highlighting}[]
\CommentTok{\# En ranger}
\BuiltInTok{set}\NormalTok{ preview\_images true}
\BuiltInTok{set}\NormalTok{ preview\_images\_method kitty}
\end{Highlighting}
\end{Shaded}

\subsection{Integración con Neovim}\label{integraciuxf3n-con-neovim}

\begin{Shaded}
\begin{Highlighting}[]
\CommentTok{{-}{-} En init.lua}
\VariableTok{vim}\OperatorTok{.}\VariableTok{opt}\OperatorTok{.}\VariableTok{termguicolors} \OperatorTok{=} \KeywordTok{true}

\CommentTok{{-}{-} Copiar a clipboard de kitty}
\VariableTok{vim}\OperatorTok{.}\VariableTok{keymap}\OperatorTok{.}\NormalTok{set}\OperatorTok{(}\StringTok{\textquotesingle{}v\textquotesingle{}}\OperatorTok{,} \StringTok{\textquotesingle{}\textless{}C{-}c\textgreater{}\textquotesingle{}}\OperatorTok{,} \StringTok{\textquotesingle{}"+y\textquotesingle{}}\OperatorTok{)}
\end{Highlighting}
\end{Shaded}

\subsection{SSH con Túneles}\label{ssh-con-tuxfaneles}

\begin{Shaded}
\begin{Highlighting}[]
\CommentTok{\# Túnel reverso para remote control}
\FunctionTok{ssh} \AttributeTok{{-}R}\NormalTok{ 52698:localhost:52698 servidor}

\CommentTok{\# O con kitten ssh (automático)}
\ExtensionTok{kitten}\NormalTok{ ssh servidor}
\end{Highlighting}
\end{Shaded}

\subsection{Diff de Directorios}\label{diff-de-directorios}

\begin{Shaded}
\begin{Highlighting}[]
\ExtensionTok{kitten}\NormalTok{ diff dir1/ dir2/}
\end{Highlighting}
\end{Shaded}

\subsection{Broadcast Input}\label{broadcast-input}

\begin{Shaded}
\begin{Highlighting}[]
\CommentTok{\# Escribir en múltiples servers}
\ExtensionTok{kitten}\NormalTok{ broadcast }\AttributeTok{{-}{-}match{-}tab} \StringTok{"title:server*"}
\end{Highlighting}
\end{Shaded}

\subsection{Quick Terminal
(Quake-style)}\label{quick-terminal-quake-style}

\begin{Shaded}
\begin{Highlighting}[]
\CommentTok{\# Crear script kitty{-}quake.sh}
\CommentTok{\#!/bin/bash}
\DataTypeTok{kitty} \DataTypeTok{{-}{-}class}\OperatorTok{=}\DataTypeTok{kitty{-}quake} \DataTypeTok{{-}{-}name} \DataTypeTok{quake} \ErrorTok{\textbackslash{}}
  \DataTypeTok{{-}o} \DataTypeTok{initial\_window\_width}\OperatorTok{=}\DecValTok{100}\DataTypeTok{c} \ErrorTok{\textbackslash{}}
  \DataTypeTok{{-}o} \DataTypeTok{initial\_window\_height}\OperatorTok{=}\DecValTok{30}\DataTypeTok{c} \ErrorTok{\textbackslash{}}
  \DataTypeTok{{-}o} \DataTypeTok{remember\_window\_size}\OperatorTok{=}\DataTypeTok{no}
\end{Highlighting}
\end{Shaded}

\textbf{Configurar en i3/sway:}

\begin{verbatim}
for_window [class="kitty-quake"] floating enable, move position center
bindsym $mod+grave exec ~/.local/bin/kitty-quake.sh
\end{verbatim}

\subsection{8. URLs Personalizadas}\label{urls-personalizadas}

\begin{Shaded}
\begin{Highlighting}[]
\CommentTok{\# Abrir issues de GitHub}
\DataTypeTok{map} \DataTypeTok{ctrl}\ErrorTok{+}\DataTypeTok{shift}\ErrorTok{+}\DataTypeTok{o} \DataTypeTok{kitten} \DataTypeTok{hints} \DataTypeTok{{-}{-}type}\OperatorTok{=}\DataTypeTok{url} \DataTypeTok{{-}{-}regex} \DataTypeTok{"}\SpecialCharTok{\textbackslash{}\textbackslash{}}\DataTypeTok{b[A{-}Z]+{-}}\SpecialCharTok{\textbackslash{}\textbackslash{}}\DataTypeTok{d+"} \DataTypeTok{{-}{-}program} \ErrorTok{@}
\end{Highlighting}
\end{Shaded}

\subsection{Layouts Personalizados}\label{layouts-personalizados}

Crear \texttt{\textasciitilde{}/.config/kitty/layouts/my\_layout.py}:

\begin{Shaded}
\begin{Highlighting}[]
\ImportTok{from}\NormalTok{ kitty.layout }\ImportTok{import}\NormalTok{ Layout}

\KeywordTok{class}\NormalTok{ MyLayout(Layout):}
\NormalTok{    name }\OperatorTok{=} \StringTok{"my\_layout"}
    
    \KeywordTok{def}\NormalTok{ do\_layout(}\VariableTok{self}\NormalTok{, windows, active\_window\_idx):}
        \CommentTok{\# Tu código de layout}
        \ControlFlowTok{pass}
\end{Highlighting}
\end{Shaded}

\subsection{Themevar para Colores
Dinámicos}\label{themevar-para-colores-dinuxe1micos}

\begin{Shaded}
\begin{Highlighting}[]
\CommentTok{\# Variables de color}
\DataTypeTok{env} \DataTypeTok{KITTY\_FG}\OperatorTok{=}\CommentTok{\#ffffff}
\DataTypeTok{env} \DataTypeTok{KITTY\_BG}\OperatorTok{=}\CommentTok{\#000000}

\CommentTok{\# Usar en config}
\DataTypeTok{foreground} \ErrorTok{$\{}\DataTypeTok{KITTY\_FG}\ErrorTok{\}}
\DataTypeTok{background} \ErrorTok{$\{}\DataTypeTok{KITTY\_BG}\ErrorTok{\}}
\end{Highlighting}
\end{Shaded}

\section{Solución de Problemas}\label{soluciuxf3n-de-problemas}

\subsection{Problemas Comunes}\label{problemas-comunes}

\subsection{Ligaduras No Funcionan}\label{ligaduras-no-funcionan}

\textbf{Verificar:}

\begin{Shaded}
\begin{Highlighting}[]
\CommentTok{\# Ver features de fuente}
\ExtensionTok{fc{-}list}\NormalTok{ : family fontformat }\KeywordTok{|} \FunctionTok{grep} \AttributeTok{{-}i} \StringTok{"fira code"}

\CommentTok{\# Verificar soporte}
\ExtensionTok{kitty}\NormalTok{ +list{-}fonts }\AttributeTok{{-}{-}psnames} \KeywordTok{|} \FunctionTok{grep} \AttributeTok{{-}i} \StringTok{"fira"}
\end{Highlighting}
\end{Shaded}

\textbf{Solución:}

\begin{Shaded}
\begin{Highlighting}[]
\CommentTok{\# Habilitar explícitamente}
\DataTypeTok{disable\_ligatures} \DataTypeTok{never}

\CommentTok{\# O instalar fuente con ligaduras}
\DataTypeTok{sudo} \DataTypeTok{apt} \DataTypeTok{install} \DataTypeTok{fonts{-}firacode}
\end{Highlighting}
\end{Shaded}

\subsection{Colores Incorrectos}\label{colores-incorrectos}

\textbf{Verificar:}

\begin{Shaded}
\begin{Highlighting}[]
\CommentTok{\# Verificar TERM}
\BuiltInTok{echo} \VariableTok{$TERM}  \CommentTok{\# Debe ser xterm{-}kitty}

\CommentTok{\# Test de colores}
\ExtensionTok{kitty}\NormalTok{ +kitten themes}
\end{Highlighting}
\end{Shaded}

\textbf{Solución:}

\begin{Shaded}
\begin{Highlighting}[]
\CommentTok{\# Forzar TERM}
\DataTypeTok{term} \DataTypeTok{xterm{-}kitty}

\CommentTok{\# O reinstalar terminfo}
\DataTypeTok{kitty} \ErrorTok{+}\DataTypeTok{kitten} \DataTypeTok{terminfo}
\end{Highlighting}
\end{Shaded}

\subsection{Imágenes No se Muestran}\label{imuxe1genes-no-se-muestran}

\textbf{Verificar:}

\begin{Shaded}
\begin{Highlighting}[]
\CommentTok{\# Test básico}
\ExtensionTok{kitten}\NormalTok{ icat /path/to/image.png}
\end{Highlighting}
\end{Shaded}

\textbf{Solución:}

\begin{Shaded}
\begin{Highlighting}[]
\CommentTok{\# Verificar OpenGL}
\ExtensionTok{glxinfo} \KeywordTok{|} \FunctionTok{grep} \StringTok{"OpenGL"}

\CommentTok{\# Actualizar drivers}
\FunctionTok{sudo}\NormalTok{ apt update }\KeywordTok{\&\&} \FunctionTok{sudo}\NormalTok{ apt upgrade}
\end{Highlighting}
\end{Shaded}

\subsection{Performance Lento}\label{performance-lento}

\textbf{Diagnóstico:}

\begin{Shaded}
\begin{Highlighting}[]
\CommentTok{\# Ver stats de rendimiento}
\ExtensionTok{kitty} \AttributeTok{{-}{-}debug{-}rendering}

\CommentTok{\# Ver FPS}
\ExtensionTok{kitty} \AttributeTok{{-}{-}debug{-}gl}
\end{Highlighting}
\end{Shaded}

\textbf{Solución:}

\begin{Shaded}
\begin{Highlighting}[]
\CommentTok{\# Ajustar delays}
\DataTypeTok{repaint\_delay} \DataTypeTok{10}
\DataTypeTok{input\_delay} \DataTypeTok{3}

\CommentTok{\# Reducir scrollback}
\DataTypeTok{scrollback\_lines} \DataTypeTok{5000}

\CommentTok{\# Deshabilitar features pesadas}
\DataTypeTok{cursor\_blink\_interval} \DataTypeTok{0}
\end{Highlighting}
\end{Shaded}

\subsection{Remote Control No
Funciona}\label{remote-control-no-funciona}

\textbf{Verificar:}

\begin{Shaded}
\begin{Highlighting}[]
\CommentTok{\# Ver si está escuchando}
\FunctionTok{netstat} \AttributeTok{{-}an} \KeywordTok{|} \FunctionTok{grep}\NormalTok{ kitty}

\CommentTok{\# O}
\ExtensionTok{lsof} \AttributeTok{{-}i}\NormalTok{ :12345}
\end{Highlighting}
\end{Shaded}

\textbf{Solución:}

\begin{Shaded}
\begin{Highlighting}[]
\CommentTok{\# Habilitar}
\DataTypeTok{allow\_remote\_control} \DataTypeTok{yes}
\DataTypeTok{listen\_on} \DataTypeTok{unix}\ErrorTok{:/}\DataTypeTok{tmp}\ErrorTok{/}\DataTypeTok{kitty}

\CommentTok{\# Verificar permisos}
\DataTypeTok{chmod} \DataTypeTok{600} \ErrorTok{/}\DataTypeTok{tmp}\ErrorTok{/}\DataTypeTok{kitty}
\end{Highlighting}
\end{Shaded}

\subsection{Fuentes Borrosas}\label{fuentes-borrosas}

\textbf{Solución:}

\begin{Shaded}
\begin{Highlighting}[]
\CommentTok{\# Ajustar composition}
\DataTypeTok{text\_composition\_strategy} \DataTypeTok{1.0} \DataTypeTok{0}

\CommentTok{\# O en macOS}
\DataTypeTok{text\_composition\_strategy} \DataTypeTok{1.7} \DataTypeTok{30}
\end{Highlighting}
\end{Shaded}

\subsection{Clipboard No Funciona}\label{clipboard-no-funciona}

\textbf{Verificar:}

\begin{Shaded}
\begin{Highlighting}[]
\CommentTok{\# Ver proveedores de clipboard}
\ExtensionTok{xclip} \AttributeTok{{-}selection}\NormalTok{ clipboard }\AttributeTok{{-}o}
\end{Highlighting}
\end{Shaded}

\textbf{Solución:}

\begin{Shaded}
\begin{Highlighting}[]
\CommentTok{\# Instalar xclip}
\FunctionTok{sudo}\NormalTok{ apt install xclip}

\CommentTok{\# O wl{-}clipboard para Wayland}
\FunctionTok{sudo}\NormalTok{ apt install wl{-}clipboard}
\end{Highlighting}
\end{Shaded}

\subsection{Debugging}\label{debugging}

\textbf{Modo verbose:}

\begin{Shaded}
\begin{Highlighting}[]
\ExtensionTok{kitty} \AttributeTok{{-}{-}debug{-}config}
\ExtensionTok{kitty} \AttributeTok{{-}{-}debug{-}rendering}
\ExtensionTok{kitty} \AttributeTok{{-}{-}debug{-}gl}
\end{Highlighting}
\end{Shaded}

\textbf{Ver logs:}

\begin{Shaded}
\begin{Highlighting}[]
\CommentTok{\# Linux}
\ExtensionTok{\textasciitilde{}/.local/state/kitty/log}

\CommentTok{\# macOS}
\ExtensionTok{\textasciitilde{}/Library/Logs/kitty}
\end{Highlighting}
\end{Shaded}

\section{Referencia Completa}\label{referencia-completa}

\subsection{Archivos de
Configuración}\label{archivos-de-configuraciuxf3n}

\begin{verbatim}
~/.config/kitty/
├── kitty.conf              # Configuración principal
├── current-theme.conf      # Tema actual (generado)
├── diff.conf               # Config para kitty diff
├── sessions/               # Archivos de sesión
│   ├── dev.kitty
│   └── personal.kitty
├── layouts/                # Layouts personalizados
│   └── my_layout.py
└── watchers.py            # Watchers de eventos
\end{verbatim}

\subsection{Variables de Entorno}\label{variables-de-entorno}

\begin{Shaded}
\begin{Highlighting}[]
\ExtensionTok{KITTY\_CONFIG\_DIRECTORY}    \CommentTok{\# Dir de config}
\ExtensionTok{KITTY\_INSTALLATION\_DIR}    \CommentTok{\# Dir de instalación}
\ExtensionTok{KITTY\_PID}                 \CommentTok{\# PID de kitty}
\ExtensionTok{KITTY\_WINDOW\_ID}          \CommentTok{\# ID de window actual}
\ExtensionTok{KITTY\_LISTEN\_ON}          \CommentTok{\# Socket de remote control}
\end{Highlighting}
\end{Shaded}

\subsection{Códigos de Escape
Útiles}\label{cuxf3digos-de-escape-uxfatiles}

\begin{Shaded}
\begin{Highlighting}[]
\CommentTok{\# Cambiar título de window}
\BuiltInTok{echo} \AttributeTok{{-}e} \StringTok{"\textbackslash{}033]2;Nuevo Título\textbackslash{}007"}

\CommentTok{\# Cambiar color}
\BuiltInTok{echo} \AttributeTok{{-}e} \StringTok{"\textbackslash{}033]11;\#ff0000\textbackslash{}007"}

\CommentTok{\# Mostrar imagen}
\BuiltInTok{echo} \AttributeTok{{-}e} \StringTok{"\textbackslash{}033]1337;File=inline=1:}\VariableTok{$(}\FunctionTok{base64}\NormalTok{ imagen.png}\VariableTok{)}\StringTok{\textbackslash{}007"}
\end{Highlighting}
\end{Shaded}

\subsection{API de Remote Control}\label{api-de-remote-control}

\textbf{Endpoints principales:}

\begin{itemize}
\tightlist
\item
  \texttt{ls} - Listar windows/tabs
\item
  \texttt{launch} - Crear window/tab
\item
  \texttt{close-window} - Cerrar window
\item
  \texttt{close-tab} - Cerrar tab
\item
  \texttt{focus-window} - Enfocar window
\item
  \texttt{focus-tab} - Enfocar tab
\item
  \texttt{get-colors} - Obtener colores
\item
  \texttt{set-colors} - Cambiar colores
\item
  \texttt{get-text} - Obtener texto
\item
  \texttt{send-text} - Enviar texto
\item
  \texttt{resize-window} - Redimensionar
\item
  \texttt{set-window-title} - Cambiar título
\item
  \texttt{goto-layout} - Cambiar layout
\end{itemize}

\subsection{Comparison Chart}\label{comparison-chart}

\begin{longtable}[]{@{}lllll@{}}
\toprule\noalign{}
Característica & Kitty & Alacritty & iTerm2 & Wezterm \\
\midrule\noalign{}
\endhead
\bottomrule\noalign{}
\endlastfoot
GPU Rendering & ✅ & ✅ & ✅ & ✅ \\
Tabs & ✅ & ❌ & ✅ & ✅ \\
Splits & ✅ & ❌ & ✅ & ✅ \\
Imágenes & ✅ & ❌ & ✅ & ✅ \\
Ligaduras & ✅ & ✅ & ✅ & ✅ \\
Remote Control & ✅ & ❌ & ✅ & ✅ \\
Extensiones & ✅ (kittens) & ❌ & ✅ & ✅ (lua) \\
Cross-platform & ✅ & ✅ & ❌ & ✅ \\
Config Language & conf & toml & GUI & lua \\
Shell Integration & ✅ & ❌ & ✅ & ✅ \\
\end{longtable}

\section{Conclusión}\label{conclusiuxf3n}

\subsection{Roadmap de Aprendizaje}\label{roadmap-de-aprendizaje}

\textbf{Semana 1 - Básicos:}

\begin{itemize}
\tightlist
\item
  Instalar y configurar
\item
  Aprender atajos principales
\item
  Configurar fuentes y colores
\end{itemize}

\textbf{Semana 2 - Intermedio:}

\begin{itemize}
\tightlist
\item
  Dominar tabs y windows
\item
  Experimentar con layouts
\item
  Configurar shell integration
\end{itemize}

\textbf{Mes 1 - Avanzado:}

\begin{itemize}
\tightlist
\item
  Usar kittens
\item
  Remote control básico
\item
  Crear sesiones
\end{itemize}

\textbf{Mes 2+ - Experto:}

\begin{itemize}
\tightlist
\item
  Crear kittens propios
\item
  Watchers
\item
  Automatización completa
\end{itemize}

\subsection{Recursos}\label{recursos}

\textbf{Documentación Oficial:}

\begin{itemize}
\tightlist
\item
  https://sw.kovidgoyal.net/kitty/
\item
  https://sw.kovidgoyal.net/kitty/conf/
\end{itemize}

\textbf{Comunidad:}

\begin{itemize}
\tightlist
\item
  GitHub: https://github.com/kovidgoyal/kitty
\item
  Discussions: https://github.com/kovidgoyal/kitty/discussions
\item
  Reddit: r/KittyTerminal
\end{itemize}

\textbf{Temas:}

\begin{itemize}
\tightlist
\item
  https://github.com/dexpota/kitty-themes
\item
  https://github.com/kdrag0n/base16-kitty
\end{itemize}

\textbf{Configs de Ejemplo:}

\begin{itemize}
\tightlist
\item
  https://github.com/search?q=kitty.conf
\end{itemize}

\subsection{Tips Finales}\label{tips-finales}

\begin{enumerate}
\def\labelenumi{\arabic{enumi}.}
\tightlist
\item
  \textbf{Empieza simple}: No copies configs complejas al inicio
\item
  \textbf{Experimenta}: Prueba diferentes layouts y kittens
\item
  \textbf{Lee la docs}: La documentación de kitty es excelente
\item
  \textbf{Usa kitten themes}: Para encontrar tu tema perfecto
\item
  \textbf{Aprovecha shell integration}: Hace la diferencia
\item
  \textbf{Explora kittens}: Son muy potentes
\item
  \textbf{Practica atajos}: La eficiencia viene de no usar el mouse
\item
  \textbf{Contribuye}: Reporta bugs, sugestiones, o crea kittens
\end{enumerate}

\section{Publicaciones Similares}\label{publicaciones-similares}

Si te interesó este artículo, te recomendamos que explores otros blogs y
recursos relacionados que pueden ampliar tus conocimientos. Aquí te dejo
algunas sugerencias:

\begin{enumerate}
\def\labelenumi{\arabic{enumi}.}
\tightlist
\item
  \href{https://chaska-x.netlify.app/operating-system/2017-05-21-comandos-de-informacion-windows/index.pdf}{\faIcon{file-pdf}}
  \href{https://chaska-x.netlify.app/operating-system/2017-05-21-comandos-de-informacion-windows}{Comandos
  De Informacion Windows}
\item
  \href{https://chaska-x.netlify.app/operating-system/2019-06-19-adb/index.pdf}{\faIcon{file-pdf}}
  \href{https://chaska-x.netlify.app/operating-system/2019-06-19-adb}{Adb}
\item
  \href{https://chaska-x.netlify.app/operating-system/2021-08-17-limpieza-y-optimizacion-de-pc/index.pdf}{\faIcon{file-pdf}}
  \href{https://chaska-x.netlify.app/operating-system/2021-08-17-limpieza-y-optimizacion-de-pc}{Limpieza
  Y Optimizacion De Pc}
\item
  \href{https://chaska-x.netlify.app/operating-system/2021-10-21-usando-apk-en-windown-11/index.pdf}{\faIcon{file-pdf}}
  \href{https://chaska-x.netlify.app/operating-system/2021-10-21-usando-apk-en-windown-11}{Usando
  Apk En Windown 11}
\item
  \href{https://chaska-x.netlify.app/operating-system/2022-05-12-gestionar-versiones-de-jdk-en-kubuntu/index.pdf}{\faIcon{file-pdf}}
  \href{https://chaska-x.netlify.app/operating-system/2022-05-12-gestionar-versiones-de-jdk-en-kubuntu}{Gestionar
  Versiones De Jdk En Kubuntu}
\item
  \href{https://chaska-x.netlify.app/operating-system/2022-07-21-instalar-tor-browser/index.pdf}{\faIcon{file-pdf}}
  \href{https://chaska-x.netlify.app/operating-system/2022-07-21-instalar-tor-browser}{Instalar
  Tor Browser}
\item
  \href{https://chaska-x.netlify.app/operating-system/2022-08-14-crear-enlaces-duros-o-hard-link-en-linux/index.pdf}{\faIcon{file-pdf}}
  \href{https://chaska-x.netlify.app/operating-system/2022-08-14-crear-enlaces-duros-o-hard-link-en-linux}{Crear
  Enlaces Duros O Hard Link En Linux}
\item
  \href{https://chaska-x.netlify.app/operating-system/2022-09-27-comandos-vim/index.pdf}{\faIcon{file-pdf}}
  \href{https://chaska-x.netlify.app/operating-system/2022-09-27-comandos-vim}{Comandos
  Vim}
\item
  \href{https://chaska-x.netlify.app/operating-system/2023-02-16-guia-de-git-y-github/index.pdf}{\faIcon{file-pdf}}
  \href{https://chaska-x.netlify.app/operating-system/2023-02-16-guia-de-git-y-github}{Guia
  De Git Y Github}
\item
  \href{https://chaska-x.netlify.app/operating-system/2023-05-02-00-primeros-pasos-en-linux/index.pdf}{\faIcon{file-pdf}}
  \href{https://chaska-x.netlify.app/operating-system/2023-05-02-00-primeros-pasos-en-linux}{00
  Primeros Pasos En Linux}
\item
  \href{https://chaska-x.netlify.app/operating-system/2023-06-17-01-introduccion-linux/index.pdf}{\faIcon{file-pdf}}
  \href{https://chaska-x.netlify.app/operating-system/2023-06-17-01-introduccion-linux}{01
  Introduccion Linux}
\item
  \href{https://chaska-x.netlify.app/operating-system/2023-06-18-02-distribuciones-linux/index.pdf}{\faIcon{file-pdf}}
  \href{https://chaska-x.netlify.app/operating-system/2023-06-18-02-distribuciones-linux}{02
  Distribuciones Linux}
\item
  \href{https://chaska-x.netlify.app/operating-system/2023-06-19-03-instalacion-linux/index.pdf}{\faIcon{file-pdf}}
  \href{https://chaska-x.netlify.app/operating-system/2023-06-19-03-instalacion-linux}{03
  Instalacion Linux}
\item
  \href{https://chaska-x.netlify.app/operating-system/2023-06-20-04-administracion-particiones-volumenes/index.pdf}{\faIcon{file-pdf}}
  \href{https://chaska-x.netlify.app/operating-system/2023-06-20-04-administracion-particiones-volumenes}{04
  Administracion Particiones Volumenes}
\item
  \href{https://chaska-x.netlify.app/operating-system/2023-07-01-atajos-de-teclado-y-comandos-para-usar-vim/index.pdf}{\faIcon{file-pdf}}
  \href{https://chaska-x.netlify.app/operating-system/2023-07-01-atajos-de-teclado-y-comandos-para-usar-vim}{Atajos
  De Teclado Y Comandos Para Usar Vim}
\item
  \href{https://chaska-x.netlify.app/operating-system/2024-07-15-instalando-specitify/index.pdf}{\faIcon{file-pdf}}
  \href{https://chaska-x.netlify.app/operating-system/2024-07-15-instalando-specitify}{Instalando
  Specitify}
\item
  \href{https://chaska-x.netlify.app/operating-system/2025-07-10-gestiona-tus-dotfiles-con-gnu-stow/index.pdf}{\faIcon{file-pdf}}
  \href{https://chaska-x.netlify.app/operating-system/2025-07-10-gestiona-tus-dotfiles-con-gnu-stow}{Gestiona
  Tus Dotfiles Con Gnu Stow}
\end{enumerate}

Esperamos que encuentres estas publicaciones igualmente interesantes y
útiles. ¡Disfruta de la lectura!






\end{document}
